% Verbatim copy of /Users/alingalatan/workspaces/grothe/handoff/manuscript/legacy/ch55-buffered-drivers.tex (file dated 2026-10-25). Only this line was added.
% Consolidated source module (formerly Chapter 55). Owner: lower-bound/10-coupled-complements.tex
\providecommand{\Sd}{\textup{[S]}}\providecommand{\cert}[1]{\texttt{#1}}

\section{Buffered drivers: two drivers at every buffer $g\ge0.164$, and any number of small drivers}\label{ch:buffered-drivers}
\chapsummary{Buffered two-driver and small-driver comparisons with retained deficits, plus joint tangents and conditional-mean-zero costs. The exact thresholds below have different hypotheses and should not be merged by taking coordinatewise maxima.}

\paragraph{Provenance and status (IDE note).} The two-driver certificate (Theorems~\ref{tdb-thm:A}, \ref{tdb-thm:B}) came from another AI working from the
state-of-the-art bundle (\code{two-driver-buffered}); ``Lemma/Theorem $n$'' and \Sd{} refer to \code{conjK\_state.pdf} in \pathref{bundles/conjK-state-20261020.zip}
(Lemma~3 visible mixture, 5 buffered tangent, 6 variance exclusion, 9 convex interpolation in $g$, 13 Bernstein deficit, Theorem~18 large buffer). The full-context GPT
audited and replayed it (two printed constants corrected below: $K\le17050$, exact $x_L$; the conclusion $2.58$ survives), then added
\S\ref{sec:buffered-review-anyn}--\S\ref{sec:tdb-downward}. Package \pathref{gpt-reviews/buffered-drivers-20261020/}. IDE replay \code{python3 -I verify\_all.py
--jobs 2} from an unrelated directory, regenerating all tables and anchors: plan A (1{,}609 boxes, worst margin $3.46\cdot10^{-8}$), plan B (455 boxes, $4.0\cdot10^{-5}$), the strengthened exact assertions (local coefficient $-2.5899$, $K=17048.5$, exact $x_L$), the trace, the seven visible constants and positive matrix, and all eight buffer-extension claims PASS (1014\,s; 17\,min wall). IDE checks \st{P}: the covariance bounds of
\eqref{eq:buffered-review-covariance}; the three final square completions in exact arithmetic ($u^\top P^{-1}u/2500=3.659,4.627,5.173$ against
$17/3-2,\,17/3-1,\,17/3-49/100$); the heat-transfer inequality $H'\ge1/H$.

\subsection{Results}

\subsubsection{Setting and what is used from the state document}
Notation is that of \Sd: $\varepsilon=1891/2500$, $p=1/\varepsilon$, $a=699/1250$, $b=179/1250$, $L=Z(a+bT)$, $T=|Z|^2$,
$H(X)=X(|X|^2-2)/\sqrt2$, $\Lambda(G)=\E|L+G|+\E|L-G|-p\E|G|^2$, $S=\sup_vF(v)$, $F(v)=2\E|L+\sqrt vU|-pv$,
$\rho_g(u)=\E|\sqrt u+gU|$, $K_g=2\rho_g'$. The family above is denoted $\mathcal F_2$; write $\theta=(c_1,c_2,\alpha_1,\alpha_2,\beta_1,\beta_2)$,
$s=c_1^2+c_2^2$ and $V=\E|G|^2=s+\sum_k(|\alpha_k|^2+|\beta_k|^2)+g^2$.

From \Sd\ we use exactly:
\begin{itemize}
\item Lemma~3 (the visible law is a Gaussian scale mixture; $F$ is strictly concave) --- for the monotonicity of $C(v)=-h''(v)$ and
of $B(v)=h'''(v)/(3C(v))$ in Theorem~\ref{tdb-thm:local};
\item Lemma~5 (concavity of $\rho_g$ and the Laplace representation of $K_g$) --- for the pointwise tangent and for the Fourier transform of $K_g$;
\item Lemma~6 (cell-tangent majorant), used with point cells, i.e.\ the pointwise tangent of the concave $\rho_g$;
\item Lemma~9 (exact convex interpolation in the buffer);
\item Lemma~13 (nonperturbative baseline deficit) --- its proof, applied to a general remote cubic (Corollary~\ref{tdb-cor:bern2});
\item Theorem~18 (large buffer, $g\ge0.477$);
\item the reduction $\Lambda(G+gU)=2\E\rho_g(|L+G|^2)-p(\E|G|^2+g^2)$ for families closed under negation (Section~3 of \Sd).
\end{itemize}
Theorems~16, 19, 20, 21 are \emph{not} used. (The checker implements an optional energy shift that would invoke Theorem~19; no box needs it.)

\subsubsection{Statements}

\begin{theorem}[Buffered two-driver family, \st{P,C}]\label{tdb-thm:A}
For every $G\in\mathcal F_2$ with $g\ge41/250$,
\[
\Lambda(G)\le S .
\]
More precisely: for $g\ge0.477$ this is Theorem~18 of \Sd; for $41/250\le g\le0.477$,
\begin{enumerate}
\item if $V\ge13/20$ then $\Lambda(G)<S-1/2000$ (Lemma~\ref{tdb-lem:trace});
\item if $0.185\le g\le0.22$ and $c_1,c_2\le0.03$ then $\Lambda(G)\le S-2.58\,(c_1^4+c_2^4)$ (Theorem~\ref{tdb-thm:local}; here $0.03$ means $\mathrm{fl}(0.03)$);
\item otherwise $\Lambda(G)\le S-3.4\cdot10^{-8}$ (Section~\ref{tdb-sec:cover}).
\end{enumerate}
Equality holds only at the Gaussian benchmark $G=\sqrt{v_*}\,U$.
\end{theorem}

\begin{theorem}[Small buffers, one large cubic, \st{P,C}]\label{tdb-thm:B}
Let $G\in\mathcal F_2$ with $\mathrm{fl}(0.1)\le g\le41/250$ (where $\mathrm{fl}(0.1)=\frac1{10}+5.6\cdot10^{-18}$) and $\max(c_1,c_2)\ge\hat c(g)$, where
\begin{center}\small
\begin{tabular}{l|ccccccc}
$g\in$ & $[0.1,0.105]$ & $[0.105,0.12]$ & $[0.12,0.135]$ & $[0.135,0.145]$ & $[0.145,0.155]$ & $[0.155,0.16]$ & $[0.16,41/250]$\\\hline
$\hat c(g)$ & $0.4$ & $0.35$ & $0.3$ & $0.25$ & $0.2$ & $0.15$ & $0.1$
\end{tabular}
\end{center}
(at a common endpoint either value applies). Then $\Lambda(G)<S-4\cdot10^{-5}$. All $\alpha_k,\beta_k\in\C$ and the other cubic coefficient are
unrestricted.
\end{theorem}
The thresholds are the binary64 numbers $\mathrm{fl}(j\cdot0.1)$ and $\mathrm{fl}(\tfrac12(\mathrm{fl}(j\cdot0.1)+\mathrm{fl}((j+1)\cdot0.1)))$ of the quadtree
grid, which exceed the decimals shown by less than $5\cdot10^{-17}$; the same convention applies to all decimal constants below
(buffer nodes, the local square side $0.03$, which is $\mathrm{fl}(0.03)<0.03$, and so on). Every check uses the exact rational value of the binary64 number.

\begin{remark}[What remains open in the two-driver family]
After Theorems~\ref{tdb-thm:A}, \ref{tdb-thm:B} and Theorems~19--21 of \Sd, the open part of $\mathcal F_2$ is
\[
0\le g<41/250,\qquad \max(c_1,c_2)<\hat c(g),
\]
with $\hat c$ the threshold of Theorem~\ref{tdb-thm:B} for $g\ge\mathrm{fl}(0.1)$ (and no restriction on $c$ for $g<\mathrm{fl}(0.1)$), intersected with $\sum_k(|\alpha_k|^2+|\beta_k|^2)<4/25$, $V<13/20$, and, for each $k$, $152c_k^2+66|\alpha_k|^2>52|\beta_k|^2$ (otherwise
Theorem~21 applies with driver $k$ as the block). Section~\ref{tdb-sec:obstruction} explains why the present method stops there.
\end{remark}

%======================================================================
\subsection{Reductions}\label{tdb-sec:red}

\begin{lemma}[Real form, \st{P}]\label{tdb-lem:real}
Rotating each $X_k$ by a phase we may take $c_k\ge0$. Rotating all coordinates by the phase of $Z$ (measure preserving), every
expectation is over $r=|Z|$ with density $2re^{-r^2}$, $\ell(r)=r(a+br^2)$, and independent standard complex $Y_k=x_ke^{i\vartheta_k}$:
\begin{multline*}
L+G\ \overset{d}{=}\ e^{i\arg Z}\Bigl(\ell(r)+\sum_k\bigl[c_k\psi(x_k)e^{i\vartheta_k}+\alpha_k(r^2-1)x_ke^{i\vartheta_k}+\beta_k\tfrac{r^2}{\sqrt2}x_ke^{-i\vartheta_k}\bigr]\Bigr)+gU,\\
\psi(x)=\frac{x(x^2-2)}{\sqrt2}.
\end{multline*}
$\mathcal F_2$ is closed under negation ($G(X)\mapsto G(-X)$) and under $X_1\leftrightarrow X_2$, so $\Lambda(G)=2\E\rho_g(|F|^2)-p(\E|G-gU|^2+g^2)$
with $F=\ell+\sum_k[\cdots]$, and it suffices to treat $c_2\le c_1$. Complex conjugation of $(Y_1,Y_2)$ maps $F$ at real parameters to $\bar F$;
hence all coefficients of the majorants below are real.
\end{lemma}
\begin{proof}
$H(e^{i\phi}Y)=e^{i\phi}H(Y)$, $(T-1)e^{i\phi}Y$, $Z^2\overline{e^{i\phi}Y}=e^{-i\phi}Z^2\bar Y$, so rotating $X_k$ multiplies $c_k,\alpha_k$ by $e^{i\phi}$
and $\beta_k$ by $e^{-i\phi}$; choose $\phi$ with $c_ke^{i\phi}\ge0$. With $Z=re^{i\omega}$ and $Y_k=e^{-i\omega}X_k$ every term acquires the
common factor $e^{i\omega}$. The rest is \Sd, Section~3.
\end{proof}

\begin{lemma}[Trace exclusion, \st{P,C}]\label{tdb-lem:trace}
For $G\in\mathcal F_2$, $\Lambda(G)\le J(V):=2\E\sqrt{T(a+bT)^2+V}-pV$. Moreover $J(13/20)<1.3713318<S-1/2000$ and $J'(V)\le V^{-1/2}-p<0$ for $V\ge13/20$.
Hence $V\ge13/20$ implies $\Lambda(G)<S-1/2000$.
\end{lemma}
\begin{proof}
$|l+w|+|l-w|\le2\sqrt{|l|^2+|w|^2}$ and $\sqrt{P+|G|^2}\le\sqrt{P+V}+(|G|^2-V)/(2\sqrt{P+V})$ with $P=T(a+bT)^2$. Given $Z$, the terms of
$G$ are centred and pairwise orthogonal ($\E[H(X)\bar X]=\E|X|^4-2=0$, $\E[H(X)X]=\E X^2=0$), so
$\E[|G|^2\mid Z]-V=\sum_k|\alpha_k|^2((T-1)^2-1)+|\beta_k|^2(T^2/2-1)$. The weight $w(T)=1/(2\sqrt{P(T)+V})$ is decreasing, and the probability
measures $(t-1)^2e^{-t}dt$, $\tfrac{t^2}2e^{-t}dt$ stochastically dominate $e^{-t}dt$ (tails $(x^2+1)e^{-x}$, $(x^2/2+x+1)e^{-x}$ versus $e^{-x}$).
Hence $\E[(T-1)^2w]\le\E w$ and $\E[\tfrac{T^2}2w]\le\E w$, and the error term is $\le0$. The value $J(13/20)$ is an upper Riemann sum on
$600{,}000$ cells of $[0,60]$ (the integrand is increasing in $t$) plus an explicit tail, in interval arithmetic (\cert{trace.cpp});
$S\ge S_{\rm lo}>1.37395352558$ (Section~\ref{tdb-sec:cover}). $J$ is concave and $J'(V)=\E(P+V)^{-1/2}-p\le V^{-1/2}-p$.
\end{proof}
This recovers, for $\mathcal F_2$, the variance exclusion of \autoref{ch:buffer-bootstrap} of \Sd; we reprove it to keep the certificate independent.

%======================================================================
\subsection{The exact Wick-anchor majorant}

Fix $g>0$ and an anchor $c^0=(c_1^0,c_2^0)$, $c_k^0\ge0$, and put $F_0=\ell(r)+\sum_kc^0_k\psi(x_k)e^{i\vartheta_k}$ (no $A/B$ part).
Order the directions as $\Phi=(H_1,H_2,A_1,A_2,B_1,B_2)$ with $H_k=\psi(x_k)e^{i\vartheta_k}$, $A_k=(r^2-1)x_ke^{i\vartheta_k}$,
$B_k=\tfrac{r^2}{\sqrt2}x_ke^{-i\vartheta_k}$ (orthonormal), so $F=F_0+\sum_j\delta_j\Phi_j$ with $\delta=\theta-\theta^0$, $\theta^0=(c^0,0,0,0,0)$.

\begin{lemma}[Global quadratic majorant, \st{P}]\label{tdb-lem:maj}
With $K=K_g(|F_0|^2)$, define the real numbers
\[
L_0=2\E\rho_g(|F_0|^2)-p(|c^0|^2+g^2),\qquad b_j=\E[K\bar F_0\Phi_j]-p\theta^0_j,\qquad M_{ij}=\E[K\bar\Phi_i\Phi_j]-p\delta_{ij}.
\]
Then for \emph{every} parameter value, with $u=\Re\delta$, $v=\Im\delta$ (note $v_H=0$),
\[
\Lambda(G)\le L_0+2b\cdot u+u^{\!\top}Mu+v^{\!\top}Mv .
\]
Let $h=c-c^0$, $w(h)=b_{ab}+M_{ab,c}h$ and $\lambda'\ge0$ with $\lambda' I-M_{ab}\succ0$. Then
\[
\sup_{\alpha,\beta}\Lambda\le L_0+2b_c\cdot h+h^{\!\top}M_{cc}h+w(h)^{\!\top}(\lambda'I-M_{ab})^{-1}w(h)+\lambda'R^2,
\]
where the supremum is over all $\alpha,\beta$ if $\lambda'=0$, and over $\sum_k(|\alpha_k|^2+|\beta_k|^2)\le R^2$ otherwise.
\end{lemma}
\begin{proof}
By Lemma~5 $\rho_g$ is concave, so pointwise $\rho_g(|F|^2)\le\rho_g(|F_0|^2)+\tfrac12K(|F|^2-|F_0|^2)$ (Lemma~6 with point cells).
Expand $|F|^2-|F_0|^2=2\Re(\bar F_0\sum\delta_j\Phi_j)+|\sum\delta_j\Phi_j|^2$ and $\E|G-gU|^2-|c^0|^2=2\theta^0\cdot u+|\delta|^2$ (orthonormality).
By Lemma~\ref{tdb-lem:real} $b$ and $M$ are real, so $\Re\sum_j b_j\delta_j=b\cdot u$ and $\delta^*M\delta=u^\top Mu+v^\top Mv$.
For the second statement complete the square in $u_{ab}$ for the matrix $M_{ab}-\lambda' I\prec0$ and use $v_{ab}^\top(M_{ab}-\lambda'I)v_{ab}\le0$,
$\lambda'(|u_{ab}|^2+|v_{ab}|^2)\le\lambda'R^2$ (Lemma~8 of \Sd).
\end{proof}
The majorant is exact at the anchor ($\delta=0$) and its gradient is exact; it overestimates the curvature in $c$ ($M_{cc}\approx0.66$ at the
benchmark versus the true value $F'(v_*)=0$). This is why the cover needs small boxes near $c=0$, and why the benchmark corner needs
Theorem~\ref{tdb-thm:local}.

\begin{lemma}[Exact box bound, \st{P}]\label{tdb-lem:box}
Let $L_0,b,M$ be given by interval enclosures and let $Q$ be the box $[c_{1a},c_{1b}]\times[c_{2a},c_{2b}]$. Let $N$ be a rational symmetric
$4\times4$ matrix and $\epsilon\ge\|(\lambda'I-M_{ab})-N\|_F$ for all matrices in the enclosure; if $N-\epsilon I\succ0$ (exact $LDL^\top$),
put $P=(N-\epsilon I)^{-1}$. Then $(\lambda'I-M_{ab})^{-1}\preceq P$, and with $w=w_0(h)+e$, $|e_i|\le\rho_i(h)$,
\[
w^\top(\lambda'I-M_{ab})^{-1}w\le w_0^\top Pw_0+2\textstyle\sum_i|(Pw_0)_i|\rho_i+\sum_{ij}|P_{ij}|\rho_i\rho_j .
\]
The resulting bound is a quadratic $q_0+2q_1\cdot h+h^\top Q_2h$ with exact rational coefficients plus an error term monotone in $|h_i|$;
its maximum over the rectangle is taken exactly (corners; critical points on edges when the edge coefficient is negative; the interior
critical point when $Q_2\prec0$).
\end{lemma}
\begin{proof}
$\lambda'I-M_{ab}=N-\Delta\succeq N-\epsilon I\succ0$ and inversion is operator monotone decreasing. The expansion of $w^\top Pw$ is exact
and the remaining steps are triangle inequalities. A quadratic on a rectangle attains its maximum at a corner, at a critical point of
its restriction to an edge, or at an interior critical point, which exists as a maximum only if $Q_2\prec0$; degenerate concave cases
attain the maximum on the boundary.
\end{proof}

%======================================================================
\subsection{Hankel factorisation of the anchor coefficients}

Write $r$ and $x$ for independent variables with density $2te^{-t^2}$ on $(0,\infty)$. The universal tables are
\[
R^{(N)}_P(k)=\E\bigl[P(r)J_N(k\ell(r))\bigr],\qquad D^{(n)}_q(s)=\E\bigl[q(x)J_n(s\psi(x))\bigr],
\]
for the visible weights
\[
(P,N)\in\{(1,0),\,(r^2-1,0),\,((r^2-1)^2,0),\,(r^4/2,0),\,(r^2/\sqrt2,2),\,((r^2-1)r^2/\sqrt2,2)\},
\]
denoted $R_1,R_A,R_{AA},R_{BB},R_B,R_{AB}$, and the driver weights
\[
(q,n)\in\{(1,0),\,(\psi^2,0),\,(x\psi,0),\,(x\psi,2),\,(x^2,0),\,(x^2,2),\,(\psi,1),\,(x,1)\},
\]
denoted $D_1,D_{HH},D_{HX},D^{(2)}_{HX},D_{XX},D^{(2)}_{XX},D^{(1)}_H,D^{(1)}_X$.
In addition, with $\kappa_Q(z)=(1-J_0(z))/z^2$, $\kappa_1(z)=J_1(z)/z$, $\kappa_R(z)=(J_1(z)/z-\tfrac12)/z^2$:
\begin{gather*}
R_Q(k)=\E[\ell^2\kappa_Q(k\ell)],\qquad Q_H(s)=\E[\psi^2\kappa_Q(s\psi)],\\
Y_H(s)=\E[\psi^2\kappa_1(s\psi)],\qquad Y_X(s)=\E[x\psi\,\kappa_1(s\psi)],\qquad Y_R(s)=\E[x\psi^3\kappa_R(s\psi)],
\end{gather*}
so that $D^{(1)}_H(s)=sY_H(s)$, $D^{(1)}_X(s)=sY_X(s)$, $Y_X(s)=s^2Y_R(s)$ (because $\E[x\psi]=0$). On the real line every table is bounded by $1$:
$|J_n|\le1$, $|\kappa_1|\le\frac12$, $0\le\kappa_Q\le\frac14$, and $\E|P(r)|\le1$, $\E|q(x)|\le1$ for all weights above (Cauchy--Schwarz with
$\E(r^2-1)^2=\E r^4/2=\E\psi^2=\E x^2=1$); moreover $|Y_R|\le\E[x|\psi|^3]/16\le\sqrt{\E x^2\psi^2\,\E\psi^4}/16=\sqrt{176}/16<0.83$.

\begin{proposition}[Hankel factorisation, \st{P}]\label{tdb-prop:hankel}
Let $g>0$, $x=gk/2$ and
\[
W_g(k)=e^{-x^2}-\sqrt\pi\,x\,\mathrm{erfc}(x),\qquad E_g(k)=e^{-g^2k^2/4}.
\]
Then, with $D(\cdot)$ evaluated at $c_1k$ for driver~1 and at $c_2k$ for driver~2 ($c=c^0$),
\begin{align*}
\E|F_0+gU|&=\int_0^\infty\Bigl[R_Q+R_1c_1^2Q_H(c_1k)+R_1D_1(c_1k)c_2^2Q_H(c_2k)+R_1D_1(c_1k)D_1(c_2k)\frac{1-E_g(k)}{k^2}\Bigr]dk,\\
b_{H_1}+pc_1&=c_1\!\int_0^\infty E_gR_1\,Y_H(c_1k)D_1(c_2k)\,dk,\qquad
b_{A_1}=c_1\!\int_0^\infty E_gR_A\,Y_X(c_1k)D_1(c_2k)\,dk,\\
b_{B_1}&=-c_1\!\int_0^\infty E_gR_B\,Y_X(c_1k)D_1(c_2k)\,dk,
\end{align*}
symmetrically for driver~2, and $M_{ij}+p\delta_{ij}=\int_0^\infty W_g(k)\,m_{ij}(k)\,dk$ with
\begin{center}\small
\begin{tabular}{llll}\toprule
same driver $k$ (other $o$) & $m_{ij}$ & cross pair & $m_{ij}$\\\midrule
$H_kH_k$ & $R_1D_{HH}(c_k\cdot)D_1(c_o\cdot)$ & $H_1H_2$ & $-R_1D^{(1)}_H(c_1\cdot)D^{(1)}_H(c_2\cdot)$\\
$H_kA_k$ & $R_AD_{HX}(c_k\cdot)D_1(c_o\cdot)$ & $H_1A_2$, $A_1H_2$ & $-R_AD^{(1)}_HD^{(1)}_X$, $-R_AD^{(1)}_XD^{(1)}_H$\\
$H_kB_k$ & $R_BD^{(2)}_{HX}(c_k\cdot)D_1(c_o\cdot)$ & $H_1B_2$, $B_1H_2$ & $+R_BD^{(1)}_HD^{(1)}_X$, $+R_BD^{(1)}_XD^{(1)}_H$\\
$A_kA_k$ & $R_{AA}D_{XX}(c_k\cdot)D_1(c_o\cdot)$ & $A_1A_2$ & $-R_{AA}D^{(1)}_XD^{(1)}_X$\\
$A_kB_k$ & $R_{AB}D^{(2)}_{XX}(c_k\cdot)D_1(c_o\cdot)$ & $A_1B_2$, $B_1A_2$ & $+R_{AB}D^{(1)}_XD^{(1)}_X$\\
$B_kB_k$ & $R_{BB}D_{XX}(c_k\cdot)D_1(c_o\cdot)$ & $B_1B_2$ & $-R_{BB}D^{(1)}_XD^{(1)}_X$\\\bottomrule
\end{tabular}
\end{center}
(in the cross column the first table is at $c_1k$, the second at $c_2k$). All integrals converge absolutely.
\end{proposition}
\begin{proof}
\emph{Transforms.} For $w\in\C=\R^2$ and $\xi=ke^{i\tau}$ write $\langle\xi,w\rangle=\Re(\bar\xi w)$. In two dimensions
$|w|=\frac1{2\pi}\int(1-\cos\langle\xi,w\rangle)|\xi|^{-3}d\xi$; for a circular law with radial characteristic function $\chi$ this gives
$\E|W|=\int_0^\infty(1-\chi(k))k^{-2}dk$. Undoing the rotation of Lemma~\ref{tdb-lem:real}, $\E|F_0+gU|=\E|L+c_1H(X_1)+c_2H(X_2)+gU|$ is the
mean modulus of a sum of four independent circular variables, so $\chi(k)=R_1(k)D_1(c_1k)D_1(c_2k)E_g(k)$, since $\E J_0(k|L|)=R_1(k)$,
$\E J_0(k|cH|)=D_1(ck)$ and $\E J_0(kg|U|)=E_g(k)$. The splitting
$1-\chi=(1-R_1)+R_1(1-D_1)+R_1D_1(1-D_1)+R_1D_1D_1(1-E_g)$ with $(1-R_1(k))/k^2=R_Q(k)$ and $(1-D_1(ck))/k^2=c^2Q_H(ck)$ is exact and free of
cancellation.

By Lemma~5, $K_g(|w|^2)=\frac1{g\sqrt\pi}\int_0^1e^{-t|w|^2/g^2}\sqrt{(1-t)/t}\,dt$; the planar transform of $e^{-t|w|^2/g^2}$ is
$(\pi g^2/t)e^{-g^2k^2/(4t)}$, hence
\[
\widehat{K_g}(k)=g\sqrt\pi\int_0^1t^{-3/2}(1-t)^{1/2}e^{-g^2k^2/(4t)}dt=\frac{2\pi}{k}\,W_g(k),
\]
the last identity by $t=1/(1+\sigma)$ and $\int_0^\infty\sigma^{1/2}e^{-x^2\sigma}(1+\sigma)^{-1}d\sigma=\sqrt\pi/x-\pi e^{x^2}\mathrm{erfc}(x)$.
In particular $0<W_g\le1$, $W_g$ is decreasing in $k$ and in $g$, $\int_0^\infty W_g=\sqrt\pi/(2g)$, and $W_g(k)\le e^{-x^2}/(2x^2)$.
For the gradient, $\varphi(w)=\E|w+gU|$ satisfies $K_g(|w|^2)\bar w=2\partial_w\varphi(w)$ and $\widehat\varphi(k)=-2\pi k^{-3}E_g(k)$.

\emph{Angular integrals.} For a radial variable $t$ and an integer $m$, $\E_\vartheta[e^{im\vartheta}e^{ikt\cos(\vartheta-\tau)}]=i^me^{im\tau}J_m(kt)$,
and $\int_0^{2\pi}e^{-iM\tau}e^{iz\cos\tau}d\tau=2\pi i^{-M}J_M(z)$. Insert the Fourier representations, integrate first over $\vartheta_1,\vartheta_2$
(each direction $\Phi_j$ carries the angular mode $+1$ for $H_k,A_k$ and $-1$ for $B_k$, the product $\bar\Phi_i\Phi_j$ the difference of
modes), then over $\tau$, which produces $J_N(k\ell)$ with $N$ the total driver mode, and finally over $r$, $x_1$, $x_2$. For example
\begin{align*}
\E[K\bar F_0B_1]&=\frac1{(2\pi)^2}\!\int_0^\infty\!\!\int_0^{2\pi}\widehat\varphi(k)\,ike^{-i\tau}\,\E_r\Bigl[\tfrac{r^2}{\sqrt2}e^{ik\ell\cos\tau}\Bigr]e^{-i\tau}\,i\,D^{(1)}_X(c_1k)D_1(c_2k)\,k\,d\tau\,dk\\
&=-\!\int_0^\infty\!\frac{E_g(k)}{k}R_B(k)D^{(1)}_X(c_1k)D_1(c_2k)\,dk,
\end{align*}
which is the stated formula since $D^{(1)}_X(s)=sY_X(s)$. The other entries are identical computations; the signs collect $i^m$ factors
and $D^{(n)}_q(-s)=(-1)^nD^{(n)}_q(s)$. Absolute convergence: $|\widehat{K_g}(k)|\le2\pi/k$, $|\widehat\varphi(k)|k\le2\pi k^{-2}E_g$, the remaining
factors are bounded by moments, and Fubini applies.
\end{proof}

\begin{remark}[Independent validation, \st{F}]
A direct product Gauss quadrature in $(r,x_1,\vartheta_1,x_2,\vartheta_2)$ (no Fourier transform) converges to the certified values as the
resolution grows: at $(c_1,c_2,g)=(0.1,0.05,0.17)$ the differences in $(L_0,b,M)$ are
$(2.7\cdot10^{-8},9.5\cdot10^{-7},8.8\cdot10^{-5})$, $(8.2\cdot10^{-10},4.2\cdot10^{-8},5.6\cdot10^{-6})$, $(1.9\cdot10^{-11},4.3\cdot10^{-9},1.1\cdot10^{-7})$,
$(2.2\cdot10^{-11},3.0\cdot10^{-10},1.9\cdot10^{-8})$ at $32^3,48^3,64^3,80\cdot72\cdot64$ nodes, and similarly at $(0.7,0.2,0.17)$.
This checks every sign in Proposition~\ref{tdb-prop:hankel}.
\end{remark}

%======================================================================
\subsection{Rigorous evaluation}\label{tdb-sec:rig}
All arithmetic is binary64 with every operation rounded outward by one ulp (\cert{ia.hpp}), compiled with \texttt{-ffp-contract=off}
and without fast-math (enforced by a preprocessor check).

\paragraph{Elementary functions and Bessel functions.} $\exp$ by $n\ln2$ reduction and Taylor degree 24 with Lagrange remainder;
$\cos,\sin$ by $2\pi$ reduction and Taylor to degree 35; $\mathrm{erfc}$ by the alternating series for $x\le2$ and by two consecutive convergents of
its Stieltjes continued fraction (which bracket the value) for $x>2$. $J_n$ ($n\le2$) by the power series for $x\le12$ (alternating tail
bounded by the first omitted term once the term ratio is below $\frac12$) and by the Hankel expansion with ten terms in each of $P,Q$
and the remainder bound of DLMF 10.17(iii) for $x>12$. The kernels $\kappa_1,\kappa_Q,\kappa_R$ use their power series for $|z|\le12$ and
the Hankel $J_0,J_1$ otherwise. Interval arguments are widened by the real Lipschitz constants $|J_n'|\le1$, $|\kappa_1'|\le\frac13\cdot\frac2\pi<\frac14$,
$|\kappa_Q'|\le\frac2{9\pi}<\frac14$, $|\kappa_R'|\le\frac{4}{90\pi}<\frac1{64}$ (from $\kappa_1(z)=\frac2\pi\int_0^1\sqrt{1-u^2}\cos(zu)du$,
$\kappa_Q(z)=\frac2\pi\int_0^1u^2k_q(zu)\frac{du}{\sqrt{1-u^2}}$, $\kappa_R(z)=-\frac2\pi\int_0^1\sqrt{1-u^2}\,u^2k_q(zu)du$, $k_q(w)=\int_0^1(1-t)\cos(wt)dt$,
$|k_q'|\le\frac16$). In complex strips: $|J_n(\zeta)|\le e^{|\Im\zeta|}$, $|\kappa_1(\zeta)|,|\kappa_Q(\zeta)|\le2e^{|\Im\zeta|}$, $|\kappa_R(\zeta)|\le\frac14e^{|\Im\zeta|}$.

\paragraph{Quadrature.} The $30$-point Gauss--Legendre rule is generated in exact rational arithmetic: each node is bracketed by two
adjacent binary64 numbers with a strict sign change of $P_{30}$, and each weight $2/((1-x^2)P_{30}'(x)^2)$ is enclosed from exact endpoint values
and a second-derivative bound; the weight sum encloses $2$ (\cert{gen\_gauss.py}). If $f$ is analytic inside the Bernstein ellipse $E_2$ of
$[a,b]$ with $|f|\le M$, then $|\int_a^bf-Q_{30}f|\le8(b-a)M4^{-30}$; $E_2$ lies in the rectangle $|\Re z-m|\le1.25h$, $|\Im z|\le0.75h$
($m,h$ the midpoint and half-length).

\paragraph{Tables.} Each table integral is truncated at $r,x\le6.5$ (tails bounded through $\Gamma(j/2+1,42.25)$ by
$8.8\cdot10^{-16}$ for the visible and $2.0\cdot10^{-14}$ for the driver family) and computed by adaptive Gauss panels with the Bernstein bound
(integrand bounds from $|e^{-z^2}|\le e^{Y^2-x_{\min}^2}$, polynomial weights, and the strip bounds of the kernels). The visible tables are
computed at the $8{,}400$ Gauss nodes of the $k$-integral ($280$ panels of length $\frac14$ on $[0,70]$), the driver tables at $25$
first-kind Chebyshev nodes on each of $560$ panels of length $\frac1{10}$ on $[0,56]$, and $Y_R$ on $30$ panels on $[0,3]$. Node values
are computed at the midpoint of the node enclosure and widened by its radius times the Lipschitz constant
$\int_0^{6.5}|P||{\rm phase}|\,2te^{-t^2}dt$ ($\le6.69$ visible, $\le61.6$ driver, $\le23.1$ for $Y_R$ including $|\kappa_R'|\le\frac1{64}$).
Strip bounds $B(\delta)=2\int_0^{6.5}|P|\,2te^{-t^2}e^{\delta|{\rm phase}|}dt$ (upper Riemann sums) are $B_V(0.1)\le11.53$, $B_D(0.05)\le26.65$,
$B_D(0.1)\le42.57$, $B_{Y_R}(0.1)\le1201.2$.

\paragraph{Chebyshev panels.} On a panel of half-length $h=\frac1{20}$ the interpolant in $25$ first-kind nodes has error at most
$4B_D(0.1)\rho^{-24}/(\rho-1)\le4.8\cdot10^{-14}$ with $\rho=2+\sqrt5$ (the ellipse with semi-minor axis $0.1$; for first-kind nodes the
aliasing identity $T_{2N\pm k}\equiv-T_k$ on the nodes gives $\|f-p\|\le2\sum_{k\ge N}|c_k|$). The interpolant is evaluated at a
binary64 point $u\in[-1,1]$ by the float recurrence $t_{m+1}=2ut_m-t_{m-1}$ with the forward bound
$|t_k-T_k(u)|\le\sum_{m<k}(k-m)E_m$, $E_m\le1.12\cdot10^{-16}(4|t_m|+|t_{m-1}|)$ (the local errors propagate through $U_{k-1-m}$, $|U_j|\le j+1$),
followed by an interval dot product. (Interval Clenshaw is not used: its wrapping effect inflates radii by $(1+\sqrt2)^{24}$.)

\paragraph{The $k$-integral.} On each $k$-panel the integrands of Proposition~\ref{tdb-prop:hankel} are bounded on the rectangle of $E_2$ by
$|R|\le B_V(0.094)$, $|D(ck)|\le B_D(\max(0.05,0.094c))$, $|E_g(k)|\le e^{-g^2\Re(k^2)/4}$, $|(1-E_g)/k^2|\le\frac{g^2}4\max(1,e^{-g^2\Re(k^2)/4})$, and
$|W_g(k)|\le(|k|/k')W_g(k')$ with $k'=\sqrt{\Re k^2}$ when $\Re k^2>0$ (from the integral representation of $W_g$), otherwise
$|W_g|\le e^{\tau^2}+\sqrt\pi|x|(1+\frac2{\sqrt\pi}|x|e^{\tau^2})$, $\tau=\Im x$. Replacing the tables by their truncations costs at most
$(t_V+2t_D)\int W_g$ (resp.\ $\int E_g$); the tail $k>70$ costs at most $\int_{70}^\infty W_g\le\frac2ge^{-x_0^2}/(4x_0^3)$,
$\int_{70}^\infty E_g\le\frac2ge^{-x_0^2}/(2x_0)$, $x_0=35g$, and for $L_0$ the exact tail $\int_{70}^\infty k^{-2}=1/70$ minus
$|\int_{70}^\infty\chi k^{-2}|\le70^{-2}\int_{70}^\infty E_g$. Typical radii of the anchor enclosures: $6\cdot10^{-11}$ ($L_0$), $10^{-12}$ ($b$),
$2\cdot10^{-11}$ ($M$).

%======================================================================
\subsection{Proof of Theorem~\ref{tdb-thm:A}}\label{tdb-sec:cover}

\subsubsection{The lower bound for $S$}
$S\ge F(g_{\rm ref}^2)$ for $g_{\rm ref}=0.20086$, and the enclosure of $L_0$ at $c=0$, $g=g_{\rm ref}$ gives
\[
S\ge S_{\rm lo}=\tfrac{3093868292922215}{2251799813685248}=1.373953525583988\ldots
\]
(consistent with $S>1.37395$ of \Sd). All cover bounds are compared with $S_{\rm lo}$.

\subsubsection{The cover}
Buffer nodes: $g_0$ the largest binary64 number below $41/250$, then $0.165,0.17,\dots,0.23$ (step $0.005$), $0.24,0.25,\dots,0.28$, $0.30,0.32,0.34$,
$0.37,0.40,0.44,0.478$ (27 nodes, 26 intervals). For each interval $[g_i,g_{i+1}]$, a quadtree refinement of the $0.1$-grid of $[0,0.8]^2$
provides boxes $Q$ with anchor $c^0$ at the centre. For each box, Lemma~\ref{tdb-lem:box} gives exact rational bounds $U_i,U_{i+1}$ for
$\sup_{c\in Q,\alpha,\beta}\Lambda$ at $g_i$ and $g_{i+1}$, and Lemma~9 of \Sd\ gives, for every fixed $G$ with $c\in Q$ and
$g=(1-t)g_i+tg_{i+1}$,
\[
\Lambda(G+gU)\le\max_{t\in[0,1]}\bigl[(1-t)U_i+tU_{i+1}+p\,t(1-t)(g_{i+1}-g_i)^2\bigr]\le S_{\rm lo},
\]
the maximum being computed exactly. Coverage is checked by replaying the quadtree with exact rational tests: a cell is either a certified
box, or lies in $c_2>c_1$ (symmetry), or satisfies $c_{1}^2+c_{2}^2\ge13/20-g_i^2$ at its lower-left corner (Lemma~\ref{tdb-lem:trace}), or lies in
$[0,0.03]^2$ for an interval inside $[0.185,0.22]$ (Theorem~\ref{tdb-thm:local}). All boxes use $\lambda'=0$: the $A/B$ block is certified negative
definite at every anchor, so no restriction on $\alpha,\beta$ and no use of Theorem~19 is needed.

\begin{center}\small
\begin{tabular}{lrr@{\qquad}lrr}\toprule
buffer interval & boxes & $\max\text{bound}-S_{\rm lo}$ & buffer interval & boxes & $\max\text{bound}-S_{\rm lo}$\\\midrule
$[g_0,0.165]$ & 85 & $-3.29\cdot10^{-5}$ & $[0.215,0.22]$ & 64 & $-5.40\cdot10^{-6}$\\
$[0.165,0.17]$ & 85 & $-4.50\cdot10^{-5}$ & $[0.22,0.225]$ & 62 & $-1.61\cdot10^{-5}$\\
$[0.17,0.175]$ & 79 & $-1.52\cdot10^{-5}$ & $[0.225,0.23]$ & 62 & $-6.99\cdot10^{-5}$\\
$[0.175,0.18]$ & 79 & $-1.01\cdot10^{-5}$ & $[0.23,0.24]$ & 56 & $-1.14\cdot10^{-5}$\\
$[0.18,0.185]$ & 103 & $-1.44\cdot10^{-6}$ & $[0.24,0.25]$ & 50 & $-1.42\cdot10^{-4}$\\
$[0.185,0.19]^\ast$ & 90 & $-3.01\cdot10^{-7}$ & $[0.25,0.26]$ & 44 & $-3.23\cdot10^{-5}$\\
$[0.19,0.195]^\ast$ & 95 & $-1.95\cdot10^{-6}$ & $[0.26,0.27]$ & 41 & $-1.36\cdot10^{-4}$\\
$[0.195,0.2]^\ast$ & 104 & $-3.46\cdot10^{-8}$ & $[0.27,0.28]$ & 37 & $-7.01\cdot10^{-5}$\\
$[0.2,0.205]^\ast$ & 90 & $-1.08\cdot10^{-7}$ & $[0.28,0.30]$ & 37 & $-6.15\cdot10^{-4}$\\
$[0.205,0.21]^\ast$ & 81 & $-1.15\cdot10^{-7}$ & $[0.30,0.32]$ & 34 & $-2.64\cdot10^{-4}$\\
$[0.21,0.215]^\ast$ & 73 & $-1.24\cdot10^{-5}$ & $[0.32,0.34]$ & 34 & $-1.71\cdot10^{-3}$\\
 & & & $[0.34,0.37]$ & 34 & $-3.48\cdot10^{-3}$\\
 & & & $[0.37,0.40]$ & 32 & $-6.85\cdot10^{-3}$\\
 & & & $[0.40,0.44]$ & 29 & $-1.11\cdot10^{-2}$\\
 & & & $[0.44,0.478]$ & 29 & $-1.85\cdot10^{-2}$\\\bottomrule
\end{tabular}

{\footnotesize $^\ast$: the square $[0,0.03]^2$ is excluded (Theorem~\ref{tdb-thm:local}). Total: 1{,}609 boxes.}
\end{center}

\emph{Binding cells.} The refinement is greedy (a box is accepted as soon as it passes), so the smallest slack, $3.5\cdot10^{-8}$, sits on
the large box $[0.10,0.15]\times[0,0.05]$ at $[0.195,0.2]$. The structurally tight cells are those just outside the local square near the
benchmark buffer: $c_1\in[0.025,0.0375]$, $c_2\le0.019$ for $0.195\le g\le0.21$ (slack $1\cdot10^{-7}$ to $2\cdot10^{-6}$, box side down to
$1/320$). There the true deficit is only $\approx20(c_1^4+c_2^4)$, while the majorant overestimates the curvature in $c$ by $\approx0.66\,|h|^2$.

\subsubsection{The local theorem at the benchmark corner}

\begin{corollary}[Lemma~13 for a remote cubic, \st{P}]\label{tdb-cor:bern2}
Let $q$ be independent of $(Z,U)$ with $\E|q|^2<\infty$. Then
$\Lambda(q+\sqrt vU)\le F(v)+F'(v)\E|q|^2-\frac{C(v)}{2B(v)^2}\E\,\mathcal E_1(B(v)|q|^2)$, $\mathcal E_1(y)=y-1+e^{-y}$.
\end{corollary}
\begin{proof}
The proof of Lemma~13 in \Sd\ shows $f_v(u)=2\E|L+\sqrt vU+\sqrt u|\le h(v)+h'(v)u-\frac C{2B^2}\mathcal E_1(Bu)$ for deterministic $u\ge0$. Since
$L+\sqrt vU$ is circular and independent of $q$, $2\E|L+\sqrt vU+q|=\E f_v(|q|^2)$; take expectations and subtract $p(v+\E|q|^2)$.
\end{proof}

\begin{theorem}[Local theorem, \st{P,C}]\label{tdb-thm:local}
For $0.185\le g\le0.22$, $c_1,c_2\in[0,0.03]$ and all $\alpha_k,\beta_k\in\C$: $\ \Lambda(G)\le S-2.58\,(c_1^4+c_2^4)$.
\end{theorem}
\begin{proof}
Put $q=c_1H(X_1)+c_2H(X_2)$, $v=g^2$, and apply Lemma~\ref{tdb-lem:maj} at the \emph{exact} anchor $c^0=c$ with only the $A/B$ directions free:
\[
\Lambda(G)\le\Lambda(q+gU)+b_{ab}^\top(-M_{ab})^{-1}b_{ab}\le\Lambda(q+gU)+\frac{|b_{ab}|^2}{\mu},
\]
provided $-M_{ab}(c,g)\succeq\mu I$ uniformly. \emph{Uniform enclosures} (\cert{local.cpp}): the integrals of Proposition~\ref{tdb-prop:hankel} are
evaluated with interval inputs over the whole box $c\in[0,0.03]^2$, $g\in[0.185,0.22]$: $W_g,E_g$ are monotone in $g$; the driver tables over
$s\in[c_ak,c_bk]$ are bounded panel by panel by the value at the centre plus the Markov bound $\sum_k|a_k|k^2$ times the half-width (plus
interpolation error); the Gauss remainders use the largest $c$ and the extreme $g$. This gives $-M_{ab}\succeq\frac1{50}I$ (exact $LDL^\top$ with
the Frobenius radius $0.0735$), and, writing $b_{A_k}=c_k^3\beta_{A_k}$, $b_{B_k}=c_k^3\beta_{B_k}$ with
$\beta_{A_k}=\int E_gR_Ak^2Y_R(c_kk)D_1(c_ok)dk$, $\beta_{B_k}=-\int E_gR_Bk^2Y_R(c_kk)D_1(c_ok)dk$,
\[
|\beta_{A_k}|\le16.611,\quad|\beta_{B_k}|\le8.066,\qquad |b_{ab}|^2/\mu\le K(c_1^6+c_2^6),\quad K=\tfrac{16.611^2+8.066^2}{1/50}\le17050 .
\]
\emph{The remote part.} By Corollary~\ref{tdb-cor:bern2} and $F(v)+F'(v)s\le F(v+s)+\frac{C(v)}2s^2\le S+\frac{C(v)}2s^2$ ($F''=-C$, $C$ decreasing),
\[
\Lambda(q+gU)\le S+\tfrac{C(v)}2s^2-\tfrac{C(v)}{2B(v)^2}\,\E\,\mathcal E_1(B(v)|q|^2).
\]
By Lemma~3 of \Sd, $h(v)=\sqrt\pi\,\E\sqrt{V+v}$, so $C(v)=\frac{\sqrt\pi}4\E(V+v)^{-3/2}$ is decreasing and
$B(v)=\frac12\E(V+v)^{-5/2}/\E(V+v)^{-3/2}$ is decreasing (moment ratios of $1/(V+v)$ are increasing). Hence on the box
$C(v)\le\bar C:=C(0.185^2)\le2.376528$, $C(v)\ge C_{\rm lo}:=C(0.22^2)\ge2.133631$, $B(v)\le B_{\rm hi}:=h'''(0.185^2)/(3C(0.185^2))\le2.700426$, with
$C(v)=\frac18\int k^2R_1e^{-vk^2/4}dk$, $h'''(v)=\frac1{32}\int k^4R_1e^{-vk^2/4}dk$ enclosed as above. The map $B\mapsto B^{-2}\E\mathcal E_1(B|q|^2)$
is decreasing ($\mathcal E_1(y)/y^2$ decreases). Averaging over the relative phase of the two cubics (Jensen, $\mathcal E_1$ convex) and using
$\mathcal E_1(x+y)=\mathcal E_1(x)+\mathcal E_1(y)+(1-e^{-x})(1-e^{-y})$,
\[
\E\,\mathcal E_1(B|q|^2)\ge\mathcal E(Bc_1^2)+\mathcal E(Bc_2^2)+(1-\phi(Bc_1^2))(1-\phi(Bc_2^2)),\qquad \phi(x)=\E e^{-xZ_3(T)},\ \mathcal E(x)=x-1+\phi(x),
\]
with $Z_3(t)=t(t-2)^2/2=|H|^2$. Since $\mathcal E(x)/x^2$ and $(1-\phi(x))/x$ are decreasing, for $x\le x_L=5626541823773821/2^{61}\approx0.0024401235\ge B_{\rm hi}\,0.03^2$:
$\mathcal E(x)\ge\gamma x^2$, $1-\phi(x)\ge\gamma'x$ with $\gamma\ge18.5823$, $\gamma'\ge0.95465$ ($\phi(x_L)$ enclosed by Gauss quadrature).
Collecting, with $c_1^2c_2^2\le\frac12(c_1^4+c_2^4)$ and $c_k^6\le0.03^2c_k^4$,
\[
\Lambda(G)-S\le(c_1^4+c_2^4)\Bigl[\tfrac{\bar C-C_{\rm lo}\gamma}{2}+\tfrac12\max\bigl(0,\bar C-\tfrac12C_{\rm lo}\gamma'^2\bigr)+K\,0.03^2\Bigr]\le-2.5898\,(c_1^4+c_2^4).
\]
All constants are checked in exact rational arithmetic (\cert{check\_local.py}).
\end{proof}

\subsubsection{Assembly}
For $g\ge0.477$ use Theorem~18 of \Sd. For $41/250\le g\le0.477$ (the cover runs up to $\mathrm{fl}(0.478)>0.477$): by symmetry $c_2\le c_1$; if $c_1^2+c_2^2+g^2\ge13/20$ then $V\ge13/20$ and
Lemma~\ref{tdb-lem:trace} applies; in the local square Theorem~\ref{tdb-thm:local} applies; every other $(c,g)$ lies in a certified box of its buffer
interval. Equality: the cover and trace bounds are strict, and in the local square $\Lambda=S$ forces $c=0$, then
$\Lambda(J+gU)\le F(g^2)$ with $J$ the $A/B$ part and the strict concavity of $F$ (Lemma~3) and of the $A/B$ quadratic force $g^2=v_*$, $J=0$.
\qed

%======================================================================
\subsection{Small buffers: proof of Theorem~\ref{tdb-thm:B}}\label{tdb-sec:B}
The same machinery applies for $g\ge0.1$ (the $k$-tail bound is effective for $35g\ge3.5$). Buffer nodes $0.1,0.105,0.11,\dots,0.16$ (step $0.005$), $0.162$, $0.164$
($\mathrm{fl}(0.164)\ge41/250$); on each interval the claimed region is $\{c_2\le c_1,\ c_1\ge\hat c\}$ minus the trace region, and cells with
$c_1\le\hat c$ are recorded as unclaimed. The exact check passes with 455 boxes, all with $\lambda'=0$:
\begin{center}\small
\begin{tabular}{lccc}\toprule
buffers & $\hat c$ & boxes & $\max\text{bound}-S_{\rm lo}$\\\midrule
$[0.100,0.105]$ & $0.40$ & $23$ & $-3.4\cdot10^{-2}$\\
$[0.105,0.110]$, $[0.110,0.115]$, $[0.115,0.120]$ & $0.35$ & $3\times28$ & $-4.2\cdot10^{-3}$, $-6.4\cdot10^{-3}$, $-8.2\cdot10^{-3}$\\
$[0.120,0.125]$, $[0.125,0.130]$, $[0.130,0.135]$ & $0.30$ & $3\times28$ & $-9.7\cdot10^{-3}$, $-1.1\cdot10^{-2}$, $-1.2\cdot10^{-2}$\\
$[0.135,0.140]$, $[0.140,0.145]$ & $0.25$ & $2\times33$ & $-2.0\cdot10^{-3}$, $-2.5\cdot10^{-3}$\\
$[0.145,0.150]$, $[0.150,0.155]$ & $0.20$ & $2\times32$ & $-3.8\cdot10^{-4}$, $-9.4\cdot10^{-4}$\\
$[0.155,0.160]$ & $0.15$ & $37$ & $-2.2\cdot10^{-4}$\\
$[0.160,0.162]$, $[0.162,0.164]$ & $0.10$ & $50+47$ & $-4.0\cdot10^{-5}$, $-6.9\cdot10^{-5}$\\\bottomrule
\end{tabular}
\end{center}
Below these thresholds the anchors near the claimed boundary either have an indefinite $A/B$ block or an exact-anchor value above $S$
(Section~\ref{tdb-sec:obstruction}); the thresholds are therefore close to optimal for this majorant.

%======================================================================
\subsection{Where and why the method stops}\label{tdb-sec:obstruction}

\paragraph{The positive $A$-curvature of the pointwise tangent \st{C at the listed nodes, F in between}.}
At $c=0$ the anchor coefficients are diagonal: $M_{A_kA_k}=\E[K_g(\ell^2)(r^2-1)^2]-p$, $M_{B_kB_k}=\E[K_g(\ell^2)r^4/2]-p$, $M_{H_kH_k}=\E K_g(\ell^2)-p$:
\begin{center}\small
\begin{tabular}{lrrrrrrrrr}\toprule
$g$ & 0.10 & 0.12 & 0.14 & 0.15 & 0.16 & 0.164 & 0.17 & 0.20 & 0.24\\\midrule
$M_{AA}(0)$ & $+0.212$ & $+0.139$ & $+0.073$ & $+0.042$ & $+0.012$ & $+0.0002$ & $-0.017$ & $-0.095$ & $-0.184$\\
$M_{BB}(0)$ & $-0.561$ & $-0.563$ & $-0.565$ & $-0.565$ & $-0.566$ & $-0.567$ & $-0.568$ & $-0.571$ & $-0.577$\\
$M_{HH}(0)$ & $1.041$ & $0.955$ & $0.875$ & $0.836$ & $0.798$ & $0.783$ & $0.762$ & $0.657$ & $0.531$\\
$S-F(g^2)$ & $1.1\!\cdot\!10^{-3}$ & $8.2\!\cdot\!10^{-4}$ & $5.2\!\cdot\!10^{-4}$ & $3.8\!\cdot\!10^{-4}$ & $2.6\!\cdot\!10^{-4}$ & $2.1\!\cdot\!10^{-4}$ & $1.5\!\cdot\!10^{-4}$ & $1.3\!\cdot\!10^{-7}$ & $3.2\!\cdot\!10^{-4}$\\\bottomrule
\end{tabular}
\end{center}
$M_{AA}(0,g)$ changes sign at $g\approx0.16407$. For smaller $g$ the majorant of Lemma~\ref{tdb-lem:maj} is unbounded in $\alpha$ near $c=0$; restricting the
$A/B$ energy to $4/25$ (Theorem~19) costs $\frac{4}{25}M_{AA}$, more than the margin $S-F(g^2)$ as soon as $M_{AA}\gtrsim1.6\cdot10^{-3}$, i.e.\ for
$g\lesssim0.1637$. The true second variation in the $A$ direction is negative (at the optimum this is Theorem~23 of \Sd; elsewhere the floating landscape
shows margins $\ge10^{-3}$ at $g=0.1$); the pointwise tangent loses the radial part of the Hessian ($\rho_g''<0$), roughly a factor~2 where
$\ell\gtrsim g$. At larger $c$ the anchors have $M_{ab}\prec0$, which gives Theorem~\ref{tdb-thm:B}.
Map of the exact-anchor value $V(c,g)=L_0+b_{ab}^\top(-M_{ab})^{-1}b_{ab}$ on a grid ($c_1,c_2$ in steps of $0.02$--$0.05$, $c_2\le c_1$;
enclosures from the engine, interpolation between grid points \st{F}): $V<S$ at every grid point with $c_1\ge0.08$ for $g=0.16$, $c_1\ge0.2$ for
$g=0.15$, $c_1\ge0.25$ for $g=0.14$, $c_1\ge0.3$ for $g=0.12$, $c_1\ge0.4$ for $g=0.1$; for smaller $c_1$ there are grid points where the $A/B$ block
is indefinite or $V>S$ (e.g.\ $V-S=+2.1\cdot10^{-4}$ at $c=(0.16,0)$, $g=0.15$).

\paragraph{What is needed below $41/250$.} The one-driver core of \Sd\ (Theorem~20, item~4) overcame the same obstruction with nontrivial cells in the
visible radius and the driver coordinate and with \emph{complex $A$ anchors}, covering $\alpha$ by boxes on which the positive curvature is paid by
the margin ($\ge10^{-3}$ at $g=0.1$). For two drivers this requires anchors with $\alpha^0_1,\alpha^0_2\ne0$; then the driver laws depend on $r$ and the
Hankel factorisation of Proposition~\ref{tdb-prop:hankel} no longer separates into one-dimensional tables (one needs two-variable driver tables
$\E[q(x)J_n(k|c\psi(x)+\alpha(r^2-1)x|)]$ in $(k,\alpha(r^2-1)/c)$). This is the concrete next step; the float landscape (margins $\ge10^{-3}$ at $g=0.1$,
$\ge1.9\cdot10^{-4}$ at $g=0.164$ for the true $\Lambda$) shows no sign of a violation.

\paragraph{No violation found.} All floating-point scans (Nelder--Mead and grid, over $(c_1,c_2,\alpha,\beta,g)$ with complex phases) stayed below $S$;
there is nothing to re-validate.

%======================================================================
\subsection{Certificate and replay}
Directory \cert{two-driver-buffered/}:
\begin{itemize}
\item \cert{src/}: \cert{ia.hpp, bessel.hpp, kernels.hpp, gauss.hpp, engine.hpp} (interval core), \cert{tables.cpp} (universal tables), \cert{anchors.cpp}
(anchor enclosures), \cert{local.cpp} (uniform enclosures, $C$, $h'''$, $\phi$), \cert{trace.cpp}; \cert{gen\_gauss.py}, \cert{boxcheck.py},
\cert{check\_cover.py}, \cert{check\_local.py} (exact rationals, \texttt{fractions} only).
\item \cert{data/}: \cert{plan\_A.json} (Theorem~\ref{tdb-thm:A}), \cert{plan\_B.json} (Theorem~\ref{tdb-thm:B}), \cert{local.jobs}. The plans list buffer nodes,
boxes and anchors only; they were produced by the floating-point planner in \cert{planner/}, which is not part of the proof.
\item \cert{verify\_all.py}: compiles, regenerates the Gauss rule and \emph{all} tables, recomputes every anchor, and runs the exact checks.
\end{itemize}
Command: \texttt{python3 -I /path/to/two-driver-buffered/verify\_all.py --out DIR --jobs 2}, from an unrelated directory. Requirements: Python~3
standard library and a C++17 compiler (no fast-math). The final clean replay took 1{,}815\,s on two cores (tables 1{,}735\,s, anchors 42\,s, exact checks 22\,s); the
log is \cert{replay.log}.

\paragraph{Independent audit.} The code was audited by a separate reviewer that had not written it, re-deriving every formula and
testing every primitive (elementary functions, Bessel functions and kernels against high-precision references; Gauss rule; strip and
ellipse bounds; the 21 entries of $M$, the 6 of $b$ and $L_0$ against a direct five-dimensional integration; the box bound against
brute force over random members of the enclosures; Lemma~9 and coverage logic; the local-theorem chain). No fatal or possibly-fatal
issue was found. It reported rounding gaps of relative size $\le3\cdot10^{-14}$, all now closed: the continued-fraction branch of
$\mathrm{erfc}$ used a rounded $x^2$; the table quadrature used a rounded panel midpoint and half-length and a non-outward ellipse rectangle;
two complex bounds for $e^{-g^2k^2/4}$ rounded the wrong way when $\Re k^2<0$; the Riemann upper sums used a rounded $-a^2$ in $e^{-a^2}$; and
the range evaluation in \cert{local.cpp} lacked a distance term for a few-ulp gap at panel ends. Duplicate-line detection and a
consistency check of the local exclusion against the local-theorem region were added. The final replay was run after these fixes;
the plans and all conclusions are unchanged (the table values moved by at most $10^{-15}$).

%======================================================================

\subsection{Unrestricted shared couplings on arbitrarily many small Wick drivers}
\label{sec:buffered-review-anyn}
\paragraph{Status and scope.}
The analytic argument below is \textbf{[P]}; the seven visible integrals
have a directed certificate \textbf{[C]}.  The retained remote deficit
is an explicit input from \autoref{ch:shared-nongaussian-new},
not a conclusion of the new numerical check.
All parameters are the fixed peak parameters
\[
 a=699/1250,\quad b=179/1250,\quad p=2500/1891,\quad
 L=Z(a+b|Z|^2),\quad H(X)=X(|X|^2-2)/\sqrt2.
\]
The variables \(Z,U,X_1,\ldots,X_N\) are independent standard proper
complex Gaussians.  The integer \(N\) is arbitrary and finite.

\begin{theorem}[Unrestricted \(A/B\), any number of small Wick drivers]
\label{thm:buffered-review-anyn}
Put \(q=\sum_{j=1}^N c_jH(X_j)\), with \(c_j\ge0\), and
\[
 G=q+gU+(|Z|^2-1)\sum_j\alpha_jX_j+
       \frac{Z^2}{\sqrt2}\sum_j\beta_j\overline{X_j}.
\]
Take \(g\ge0\), and write \(e=\sum_jc_j^2\).
If
\[
 \sum_jc_j^2\le\frac{16}{2500},\qquad
 \max_j c_j\le\frac1{50},\qquad
 \frac7{200}\le g^2\le\frac9{200},
\]
then, for every \(\alpha,\beta\in\mathbb C^N\), without an amplitude
or phase restriction,
\[
 \boxed{\quad
 \Lambda(G)\le S-2\sum_jc_j^4-
                  \frac{\|\alpha\|_2^2+\|\beta\|_2^2}{2000}.
 \quad}
\]
More generally, the variance allowance can be increased to \(e\le1/100\),
with the quartic coefficient \(2\) replaced by \(49/100\).
On \(e\le9/1000\) the quartic coefficient \(1\) is valid.
The conclusion also covers \(c=0\). It is not a theorem for an arbitrary
internally coupled cubic tensor.
\end{theorem}

\paragraph{What changes.}
The earlier theorem in \autoref{ch:shared-nongaussian-new} imposed
\(\|\alpha\|^2+\|\beta\|^2\le10^{-6}\) on this same
remote-variance, coefficient and buffer range.  The present argument
removes that restriction and improves its quartic deficit from
\(3/4\) to \(2\).  It uses a global tangent inequality, rather than a
finite Taylor expansion in the couplings.

\subsubsection{A dimension-free Fourier covariance estimate \textbf{[P]}}
Write \(e=\sum c_j^2\), and, for \(k\ge0\), set
\[
 D(k)=\mathbb E[XX^*\cos(k\Re q)],\qquad
 C(k)=\mathbb E[XX^\top\cos(k\Re q)].
\]
Then
\begin{equation}
 0\preceq I-D(k)\preceq k^2eI,\qquad
 \|C(k)\|_{\mathrm{op}}\le k^2e.
 \label{eq:buffered-review-covariance}
\end{equation}
Indeed, first-chaos orthogonality of each Wick monomial gives the
exact finite-sum identity
\[
 \mathbb E[|q|^2XX^*]
       =\operatorname{diag}(e+3c_1^2,\ldots,e+3c_N^2).
\]
Under the common phase rotation \(X\mapsto e^{i\theta}X\), the
terms involving \(q^2\) and \(\bar q^2\) in
\(\mathbb E[XX^*(\Re q)^2]\) vanish. Thus this matrix is half
the displayed diagonal matrix. Applying
\(0\le1-\cos t\le t^2/2\) proves the first assertion.
For the second, \(\mathbb E XX^\top=0\), so \(C=-\mathbb E[XX^\top w]\),
where \(w=1-\cos(k\Re q)\ge0\).
Weighted Cauchy--Schwarz bounds its operator norm by that of
\(\mathbb E[XX^*w]=I-D\).  No entrywise summation depending on \(N\)
is used.

\subsubsection{Visible averaging before estimating the matrix \textbf{[P,C]}}
After the joint phase rotation of \(Z\) and the remote variables,
put \(\ell(T)=\sqrt T(a+bT)\), and define
\[
 R_{AA}(k)=\mathbb E[(T-1)^2J_0(k\ell(T))],\quad
 R_{BB}(k)=\mathbb E[T^2J_0(k\ell(T))/2],
\]
\[
 R_{AB}(k)=\mathbb E[T(T-1)J_2(k\ell(T))/\sqrt2].
\]
Let
\[
 W_g(k)=e^{-g^2k^2/4}
       -\sqrt\pi\,\frac{gk}{2}\operatorname{erfc}(gk/2).
\]
For \(K_g(u)=2\partial_u\mathbb E|\sqrt u+gU|\), the global
tangent inequality at \(\ell+q\) gives
\begin{equation}
 \Lambda(G)\le\Lambda(q+gU)
       +2\Re(b_A\cdot\alpha+b_B\cdot\beta)
       +(\alpha,\beta)^*M(q,g)(\alpha,\beta).
 \label{eq:buffered-review-global-tangent}
\end{equation}
The matrix is interpreted with the harmless conjugations in the
\(B\) coordinates dictated by the displayed model. The bound below
is invariant under those conjugations.
The Fourier representation of \(K_g\), followed by the visible
angular average, gives diagonal blocks
\(\int W_gR_{AA}D\,dk-pI\) and the transpose/conjugate
of \(\int W_gR_{BB}D\,dk-pI\).
Its cross block is \(\int W_gR_{AB}C\,dk\), up to a sign and
conjugation.  These formulas follow from
\(\mathbb E_\theta e^{im\theta+it\cos\theta}=i^mJ_m(t)\).
In particular, at \(q=0\) the cross block is zero.

For \(g_0=37/200\), define
\[
 N_A=\int_0^\infty k^2W_{g_0}(R_{AA})_-\,dk,\quad
 N_B=\int_0^\infty k^2W_{g_0}(R_{BB})_-\,dk,\quad
 N_{AB}=\int_0^\infty k^2W_{g_0}|R_{AB}|\,dk.
\]
Equation~\eqref{eq:buffered-review-covariance} proves, with
\(A=\|\alpha\|_2,\ B=\|\beta\|_2\),
\[
 (\alpha,\beta)^*[M(q,g)-M(0,g)](\alpha,\beta)
 \le e(N_AA^2+2N_{AB}AB+N_BB^2).
\]
Here \(W_g\le W_{g_0}\) for \(g\ge g_0\).
Also \(M(0,g)\preceq M(0,g_0)\): the positive representation
\[
 K_g(u)=\frac1{\sqrt\pi}\int_0^{1/g^2}
             s^{-1/2}\sqrt{1-g^2s}\,e^{-us}\,ds
\]
shows that \(K_g(u)\) decreases with \(g\), pointwise in \(u\).

The directed checker proves
\[
 M_{AA}(0,g_0)<-\frac{57}{1000},\qquad
 M_{BB}(0,g_0)<-\frac{569}{1000},
\]
\[
 N_A<\frac{2097}{1000},\qquad N_B<\frac{2939}{1000},\qquad
 N_{AB}<\frac{2791}{500}.
 \tag{*}
\]
For clarity, it does not apply an analytic quadrature remainder to
the nonsmooth functions \(R_-\) or \(|R|\).
With \(d\mu=k^2W_{g_0}(k)\,dk\) and
\(\mu(\mathbb R_+)=\sqrt\pi/(2g_0^3)\), it uses
\[
 \int R_-\,d\mu\le
 \frac{\sqrt{\mu(\mathbb R_+)\int R^2\,d\mu}-\int R\,d\mu}{2},
 \qquad
 \int |R|\,d\mu\le
 \sqrt{\mu(\mathbb R_+)\int R^2\,d\mu}.
\]
Only the analytic integrands \(R\) and \(R^2\) are quadrature inputs.

\subsubsection{Contents and the exact square completion \textbf{[P,C]}}
Set \(k_r=\sum_jc_j^r\).  The cancellation
\(\mathbb E[\overline X H(X)]=0\), the Bessel remainder
\(|J_1(t)-t/2|\le |t|^3/16\), and
\(\mathbb E|XH(X)^3|\le\sqrt{176}\) give
\[
 \|b_A\|_2\le \frac{9283}{500}\sqrt{k_6},\qquad
 \|b_B\|_2\le \frac{1127}{125}\sqrt{k_6}.
 \tag{**}
\]
To check these constants, use the proved signs
\(R_A=\mathbb E[(T-1)J_0(k\ell)]\le0\) and
\(R_B=\mathbb E[TJ_2(k\ell)/\sqrt2]\ge0\).
The respective upper integrals are
\[
 \frac{\sqrt{176}}{16}\int_0^\infty k^2e^{-g_0^2k^2/4}
                  (-R_A(k))\,dk,\quad
 \frac{\sqrt{176}}{16}\int_0^\infty k^2e^{-g_0^2k^2/4}
                  R_B(k)\,dk.
\]
All other Wick factors contribute characteristic functions of modulus
at most one. Hence these bounds are independent of \(N\).

Put \(e_0=16/2500,\ \tau=1/2000,\ u=(9283/500,1127/125)^\top\), and
\[
 P=
 \begin{pmatrix}57/1000&0\\0&569/1000\end{pmatrix}
 -e_0
 \begin{pmatrix}2097/1000&2791/500\\2791/500&2939/1000\end{pmatrix}
 -\tau I .
\]
This rational matrix is positive definite. The quadratic bound
above is valid for the nonnegative vector \((A,B)\); replacing \(e\)
by \(e_0\) only increases its upper bound on that quadrant.
The retained remainder in~\eqref{eq:buffered-review-global-tangent}
is at most
\[
 2\sqrt{k_6}\,u^\top(A,B)^\top-(A,B)P(A,B)^\top
 \le k_6u^\top P^{-1}u.
\]
The previously certified remote Fourier deficit is
\[
 \Lambda(q+gU)\le S-\frac{17}{3}k_4
\]
on exactly the coefficient, variance and buffer range of the theorem.
Finally \(k_6\le k_4/2500\), and exact rational arithmetic gives
\[
 \frac{17}{3}-\frac{u^\top P^{-1}u}{2500}
 =\frac{35141036731739}{17503094800500}>2.
\]
This proves Theorem~\ref{thm:buffered-review-anyn}.
The remote deficit actually holds up to \(e=1/100\), as proved in
the ``sharper buffered Fourier deficit'' subsection of the coherent
Wick chapter. Replacing \(e_0\) in the same rational matrix by
\(9/1000\) or \(1/100\) gives remaining quartic coefficients greater
than \(1.0399\) or \(0.49405\), respectively. The exact checker
compares these fractions with \(1\) and \(49/100\).

\subsubsection{An unrestricted-amplitude reduction for arbitrary coupled cubics}
\label{sec:buffered-review-general}
The matrix part of the argument is not restricted to independent
Wick sums.  If \(q\) is any pure remote \(H_{2,1}\) polynomial
with variance \(e\), and \(\ell\) is any unit first-chaos direction,
then
\[
 \mathbb E[|\ell|^2|q|^2]
  =e+\|D_\ell q\|_2^2+\|\bar D_\ell q\|_2^2\le4e.
\]
For a direct finite-coefficient proof, rotate \(\ell\) to \(X_1\)
and expand \(q=\sum_{|\mu|=2,|\nu|=1}a_{\mu,\nu}H_{\mu,\nu}\).
Multiplication by \(X_1\) splits into orthogonal degrees four
and two. The squared norm of the creation term is
\(\sum |a_{\mu,\nu}|^2\mu!\nu!(\mu_1+1)\), and that of
the annihilation term is
\(\sum |a_{\mu,\nu}|^2\mu!\nu!\nu_1\).
Their sum is at most \(4e\). Thus
the operator inequalities~\eqref{eq:buffered-review-covariance}
hold for arbitrary pure \(H_{2,1}\) cubics.

\begin{proposition}[All amplitudes, arbitrary shared cubic, \textbf{[P,C]}]
Let \(q\in H_{2,1}(\mathbb C^n)\) be independent of the visible \(Z\),
and let \(U\) be fresh. Assume \(e=\mathbb E|q|^2\le16/2500\)
and \(g\ge37/200\). For arbitrary coefficient vectors
\(\alpha,\beta\in\mathbb C^n\), put
\[
 D=(|Z|^2-1)\sum_j\alpha_jX_j+
       Z^2\sum_j\beta_j\bar X_j/\sqrt2 .
\]
Use the actual buffered baseline in the first variation:
\[
 {\cal C}_q(D)=\mathbb E\!\left[
   K_g(|L+q|^2)\operatorname{Re}(\overline{L+q}\,D)\right].
\]
Then
\[
 \boxed{\quad
 \Lambda(q+gU+D)\le
 \Lambda(q+gU)+2{\cal C}_q(D)
       -\frac{\|\alpha\|_2^2+\|\beta\|_2^2}{25}.
 \quad}
\]
There is no independence assumption between \(q\) and the linear
drivers in \(D\).
With the larger allowance \(e\le1/100\), the same conclusion holds
with \(1/25\) replaced by \(1/40\).
\end{proposition}
\begin{proof}
The terms in \(q+D\) change sign under \(X\mapsto-X\), with \(Z\)
fixed; also replace \(U\) by \(-U\). Thus the two norm terms agree.
Alternatively apply the tangent argument to both signs separately.
Apply the preceding covariance estimate and (*).
For the nonnegative norm vector \((A,B)\), the residual matrix is
\[
 P_0=
 \begin{pmatrix}57/1000&0\\0&569/1000\end{pmatrix}
 -\frac{16}{2500}
 \begin{pmatrix}2097/1000&2791/500\\2791/500&2939/1000\end{pmatrix}.
\]
The diagonal entries of \(P_0-I/25\) are positive, and its
determinant is exactly \(214770369/390625000000>0\).
Thus \(P_0\succeq I/25\), proving the claim.
For the larger variance allowance, replace \(16/2500\) by \(1/100\)
and check the two principal minors of \(P_0-I/40\); both are positive
in exact rational arithmetic.
\end{proof}
\paragraph{What this does not prove.}
The proposition removes the finite-amplitude difficulty on this
small-variance range, but the actual contents still have to be paid
by the remote deficit for a general coupled tensor.
It does not assert a Gaussian comparison for arbitrary \(q\), or
reuse the Wick-specific bound (**) for such \(q\).

\paragraph{Paid coupled complements, any number of core coordinates.}
Let \(Q\) denote the nonbuffered expression in
Theorem~\ref{thm:buffered-review-anyn}, let
\(\mathcal A=\sigma(Z,X_1,\ldots,X_N)\), and suppose that \(J\in L^2\) and
\(\mathbb E[J\mid\mathcal A]=0\), with \(U\) independent of
\((\mathcal A,J)\). The conditional-complement lemma below gives
an additional cost of less than \((7/2)\|J\|_2^2\).
Consequently, in the \(e\le1/100\) region, exact K still holds for
the genuinely coupled residual \(Q+J+gU\) whenever
\[
 \|J\|_2^2\le\frac7{50}k_4+
          \frac{\|\alpha\|_2^2+\|\beta\|_2^2}{7000}.
\]
In the \(e\le16/2500\) region, replace \(7/50\) by \(4/7\).
The complement need not be independent of the core and may contain
mixed core/external-coordinate cubic monomials.

\paragraph{Replay and limitations.}
\texttt{visible\_constants.cpp} reads the freshly generated visible
Hankel table, with its Gauss-node enclosures, analytic remainders and
tails, and computes the seven directed bounds in (*) and (**).
\texttt{check\_constants.py} verifies the comparisons and the final
rational matrix calculation.
The new argument uses neither the two-driver cover nor an assumption
that the actual residual is a Gaussian variance mixture.
It does use the established remote Fourier deficit and visible-kernel
signs. Its scope does not include arbitrary mixed cubic tensors or
unrestricted small-buffer amplitudes. The later even/diagonal extensions
and their exceptions are in \autoref{ch:all-sector-audit-extension}; a
complete row cover is still required.

\subsection{Retaining the local shared-sector deficit and paying for a coupled complement}
\label{sec:buffered-two-driver-retained}

Write
\[
 Q=\sum_{j=1}^2\left[c_jH(X_j)+\alpha_j(T-1)X_j+
                  \beta_jZ^2\bar X_j/\sqrt2\right],\qquad
 k_4=c_1^4+c_2^4,\quad
 e_{AB}=\sum_{j=1}^2(|\alpha_j|^2+|\beta_j|^2).
\]
Here \(Z,X_1,X_2,U\) are independent standard proper complex
Gaussians, \(T=|Z|^2\), and the buffer \(gU\) is fresh.
In this section the local endpoints \(c_L,g_-,g_+\) mean exactly
the binary64 numbers
\[
 c_L=\mathrm{fl}(0.03),\qquad
 g_-=\mathrm{fl}(0.185),\qquad
 g_+=\mathrm{fl}(0.22).
\]
In particular \(c_L<3/100\) and \(g_->1849/10000\).
The interval computations are the same ones as the incoming local
theorem; all subsequent inequalities below are rational arithmetic.

\begin{theorem}[Retained local \(A/B\) penalty]
\label{thm:buffered-two-driver-retained}
For \(0\le c_1,c_2\le c_L\), \(g_-\le g\le g_+\), and every
\(\alpha_j,\beta_j\in\mathbb C\),
\[
 \boxed{\Lambda(Q+gU)\le S-2k_4-\frac{e_{AB}}{2000}.}
\]
The alternative bound
\[
 \Lambda(Q+gU)\le S-\frac52k_4-\frac{e_{AB}}{10000}
\]
also holds.
\end{theorem}
\begin{proof}
Set \(q=c_1H(X_1)+c_2H(X_2)\).
The local computation certifies \(-M_{AB}\succeq\mu I\),
\(\mu=1/50\), and
\[
 \|b_{AB}\|^2\le
 (16.611^2+8.066^2)(c_1^6+c_2^6).
\]
These are conservative rational upper bounds on the exact enclosures.
The baseline part of the incoming proof gives
\[
 \Lambda(q+gU)\le S-Dk_4,
\]
where the following rational number is admissible:
\[
 D=-\frac{2.376528-2.133631(18.5823)}2
   -\frac12\max\left\{0,\,
        2.376528-\frac{2.133631(0.95465)^2}{2}\right\}
  =\frac{28693571761768519}{1600000000000000}.
\]
For \(z=(\alpha_1,\alpha_2,\beta_1,\beta_2)=u+iv\), the simultaneous
complex-sector tangent, at the exact cubic anchor, is
\[
 \Lambda(Q+gU)
 \le\Lambda(q+gU)+2b_{AB}^{\mathsf T}u
                   -\mu(\|u\|^2+\|v\|^2).
\]
For \(0<\tau<\mu\), completing only \(\mu-\tau\) of this square gives
\[
 2b_{AB}^{\mathsf T}u-\mu(\|u\|^2+\|v\|^2)
 \le\frac{\|b_{AB}\|^2}{\mu-\tau}-\tau e_{AB}.
\]
Since \(c_j^6\le(3/100)^2c_j^4\),
\[
 \Lambda(Q+gU)\le
 S-\left[D-
 \frac{(16.611^2+8.066^2)(3/100)^2}{\mu-\tau}\right]k_4-\tau e_{AB}.
\]
For \(\tau=1/2000\), the bracket is exactly
\[
 \frac{45670182982990747}{20800000000000000}>2.19568>2.
\]
For \(\tau=1/10000\), it is
\[
 \frac{799827031791935281}{318400000000000000}>2.51201>\frac52.
\]
Thus both conclusions retain positive slack even after the
deliberately conservative scalar rounding.
\end{proof}

\begin{lemma}[Conditional mean-zero complement cost]
\label{lem:buffered-conditional-complement}
Let \(\mathcal A\) be a sigma algebra with respect to which \(L,Q\)
are measurable. Let \(J\in L^2(\mathbb C)\) satisfy
\(\mathbb E[J\mid\mathcal A]=0\), and let \(U\) be a standard proper
complex Gaussian independent of the joint variables
\((\mathcal A,J)\). For \(g>0\),
\[
 \boxed{\Lambda(Q+J+gU)
 \le \Lambda(Q+gU)
       +\left(\frac{\sqrt\pi}{2g}-p\right)\|J\|_2^2.}
\]
No independence of \(J\) from \(Q\) is assumed.
\end{lemma}
\begin{proof}
Put \(\Psi_g(w)=\mathbb E_U|w+gU|=\rho_g(|w|^2)\).
The tangent of the concave function \(\rho_g\), with
\(K_g=2\rho_g'\), gives for deterministic \(w,j\in\mathbb C\)
\[
 \Psi_g(w+j)\le
 \Psi_g(w)+K_g(|w|^2)\Re(\bar wj)
                  +\tfrac12K_g(|w|^2)|j|^2.
\]
Apply this once with \((w,j)=(L+Q,J)\), and once with
\((w,j)=(L-Q,-J)\). The linear terms vanish after conditioning on
\(\mathcal A\). The two quadratic terms together are at most
\[
 K_g(0)\mathbb E|J|^2
 =\frac{\sqrt\pi}{2g}\|J\|_2^2,
\]
because \(K_g\) is decreasing and its Laplace representation gives
\[
 K_g(0)=\frac1{g\sqrt\pi}
       \int_0^1t^{-1/2}(1-t)^{1/2}\,dt
       =\frac{\sqrt\pi}{2g}.
\]
Conditional centering also gives
\(\mathbb E(\bar QJ)=0\), while the fresh buffer is orthogonal to
both \(Q\) and \(J\). Consequently the variance price increases by
exactly \(p\|J\|_2^2\). This proves the inequality, directly for both
signs in \(\Lambda\), without a symmetry assumption on \(J\).
\end{proof}

\begin{corollary}[A coupled complement around the two-driver local core]
\label{cor:buffered-two-driver-complement}
In the region of Theorem~\ref{thm:buffered-two-driver-retained},
let \(\mathcal A=\sigma(Z,X_1,X_2)\), and let \(J\) satisfy the
conditions of Lemma~\ref{lem:buffered-conditional-complement}.
Then
\[
 \Lambda(Q+J+gU)
 \le S-2k_4-\frac{e_{AB}}{2000}+\frac72\|J\|_2^2.
\]
In particular exact K holds whenever
\[
 \boxed{\|J\|_2^2\le\frac47k_4+\frac{e_{AB}}{7000}.}
\]
This applies to arbitrary-dimensional cubic complements containing
mixed core/external-coordinate terms, provided their conditional
expectation onto the core vanishes.
\end{corollary}
\begin{proof}
Using \(\pi<22/7<(1773/1000)^2\), \(p=2500/1891\), and
\(g>1849/10000\), the complement cost is less than
\[
 \frac{1773/1000}{2(1849/10000)}-\frac{2500}{1891}
 =\frac{12141215}{3496459}<\frac72.
\]
Combine the preceding two results.
\end{proof}

\begin{corollary}[Paid complement in the small-buffer large-cubic region]
Assume the conditions of incoming Theorem T2, so that
\(\Lambda(Q+gU)<S-4\cdot10^{-5}\), with \(g\ge\mathrm{fl}(0.1)\).
For a complement as in Lemma~\ref{lem:buffered-conditional-complement},
\[
 \|J\|_2^2\le\frac1{200000}
 \quad\Longrightarrow\quad \Lambda(Q+J+gU)<S.
\]
\end{corollary}
\begin{proof}
The same rational comparison gives
\(\sqrt\pi/(2g)-p<8\).
The added cost is therefore less than
\(8/200000=4\cdot10^{-5}\), which is paid by the certified core
deficit.
\end{proof}

\subsection{Downward transfer in the buffer and new two-driver regions}\label{sec:tdb-downward}

This note derives a further consequence of the new two-driver anchor
certificate. It does not assume any independence between a cubic and its
own A/B terms. The only new variable is the same fresh proper Gaussian
buffer already present in the family.

Let \(Q\) be square integrable and let \(U\) be a fresh proper standard
complex Gaussian independent of \((L,Q)\). Write
\[
 e=\mathbb E|Q|^2,\quad
 H(v)=\mathbb E|L+Q+\sqrt vU|
      +\mathbb E|L-Q+\sqrt vU|,\qquad
 \Lambda_v=H(v)-p(e+v),\quad p=2500/1891.
\]
The function \(H\) increases, by conditioning on the added Gaussian
increment and Jensen's inequality. In particular, for \(0\le v\le w\),
\[
 \Lambda_v\le\Lambda_w+p(w-v).
\]
The following improved transfer spends less of a retained deficit:
\[
 \boxed{\quad
 \Lambda_w\le M\le11/8,\quad
 e+v<13/20,\quad
 0\le v\le w\le(41/250+10^{-15})^2
 \quad\Longrightarrow\quad
 \Lambda_v\le M+(p-11/25)(w-v).
 \quad}                                                     \tag{HT}
\]

Indeed, Gaussian convolution with variance parameter \(t\) has generator
\(\tfrac14\Delta\), and \(\Delta|z|=1/|z|\) in the plane. Thus, for \(t>0\),
\[
 H'(t)=\frac14\left(
  \mathbb E|L+Q+\sqrt tU|^{-1}
 +\mathbb E|L-Q+\sqrt tU|^{-1}\right)\ge\frac1{H(t)}.
\]
The inequality follows from Jensen for each reciprocal and
\((1/x+1/y)/4\ge1/(x+y)\).
The distributional identity \(\Delta|z|=1/|z|\), whose right-hand
side is locally integrable in the plane, justifies the derivative
after Gaussian convolution at \(t>0\); Tonelli permits expectation.
Continuity at \(t=0\) follows from the Lipschitz property of the norm.
Monotonicity gives, for \(v<t<w\),
\[
 H(t)\le H(w)\le M+p(e+w)
 < \frac{11}8+p\left(\frac{13}{20}
              +(41/250+10^{-15})^2\right)<\frac{25}{11}.
\]
The last inequality is checked with exact fractions. Hence
\(H'(t)>11/25\); integration and continuity at \(t=0\) prove (HT).
No rotational symmetry or independence of \(L,Q\) is needed.

\paragraph{The trace shell must be covered at the lower buffer.}
An anchor cover at \(w\) normally discards
\(c_1^2+c_2^2+w\ge13/20\). This exclusion cannot simply be
carried down to \(v\). The additional checker instead covers every
rational cell with
\[
 c_1^2+c_2^2+v_{\min}<13/20,
\]
using the same globally valid anchor quadratics at \(w\), and uses
the trace theorem only for the \emph{original} law when its total
variance is at least \(13/20\). It also extends the coordinate
square beyond \(0.8\), because at zero buffer the trace radius is
\(\sqrt{13/20}>0.8\).

\paragraph{Certificate design.}
The new checker reads regenerated interval anchor data, never the
decimal margins in the author log. Its exact rational grid covers
\([c_{\min},0.9]\times[0,0.9]\), after symmetry, and recursively
bisects a cell only when needed. Every accepted cell satisfies
\[
 U_{\rm anchor}+(p-11/25)(g_{\rm anchor}^2-g_{\min}^2)
 \le S_{\rm lo}-d.
\]
The quadratic elimination is the original exact rational
\texttt{box\_bound} with no A/B energy shift. All A and B coefficients
therefore remain unrestricted and complex.

\paragraph{Certified extension theorem [P,C].}
For the incoming two-driver family, with every complex A/B coefficient
unrestricted, the following regions satisfy \(\Lambda(G)\le S\):
\[
\begin{array}{c|c|c}
 \text{condition on }\max(c_1,c_2)&
 \text{allowed buffers }g&\text{new compact cells}\\ \hline
 \ge63/200=0.315&[0,\infty)&29\\
 \ge3/10&[57/1000,\infty)&29\\
 \ge1/4&[31/250,\infty)&32\\
 \ge1/5&[143/1000,\infty)&34\\
 \ge3/20&[771/5000,\infty)&36\\
 \ge1/10&[3997/25000,\infty)&46\\
 \text{unrestricted}&[16389/100000,\infty)&81
\end{array}
\]
In addition, without any cubic-amplitude restriction,
\[
 \boxed{\quad g\ge8179/50000=0.16358
   \quad\Longrightarrow\quad \Lambda(G)<S+89/10^6.\quad}
\]
The latter check uses \(81\) cells. All displayed thresholds are exact
rationals, not binary64 conventions.

The anchor buffers for the first six rows are respectively the exact
binary64 values of \(0.13,0.13,0.135,0.145,0.155,0.16\).
For the unrestricted rows the anchor is the incoming first A node,
\(\mathrm{nextdown}(\mathrm{fl}(0.164))
 =0\mathrm{x}1.4fdf3b645a1cap{-}3<41/250\).
The new covers have been replayed against the independently regenerated
anchor matrices. Their respective retained compact deficits, rounded
here only for reporting, are
\[
 4.17526\cdot10^{-6},\quad9.38731\cdot10^{-5},\quad
 7.05794\cdot10^{-5},\quad1.07094\cdot10^{-4},\quad
 4.99012\cdot10^{-6},\quad6.20256\cdot10^{-6},\quad
 1.07457\cdot10^{-6}.
\]
Every exact-K row adopts only the strict compact deficit \(10^{-6}\).
The weak row has exact certified excess less than
\(8.84676\cdot10^{-5}<89/10^6\). The trace branch of the
\emph{original} law has deficit \(1/2000\).

\paragraph{Joining the full buffer half-lines.}
For each of the first six rows, the new argument covers the interval
from the displayed lower buffer through its anchor buffer. The incoming
T2 intervals then cover continuously up to \(41/250\), and T1 covers
\([41/250,\infty)\).
Some T2 cubic thresholds exceed the displayed exact decimal by less
than \(10^{-15}\). If necessary raise one cubic coefficient by
\(\delta\le10^{-15}\). On the original compact branch \(c<1\), and
the triangle inequality, \(\mathbb E|H|\le1\), and the orthogonal
variance price imply
\[
 \Lambda(c)\le\Lambda(c+\delta)
          +2\delta+p(2c\delta+\delta^2)
 <\Lambda(c+\delta)+6\cdot10^{-15}.
\]
This is absorbed by T2's \(4\cdot10^{-5}\) deficit. This argument is
used only to join source thresholds, not to accept a new grid cell.

For the unrestricted rows, the incoming first A interval starts at the
same binary64 node used by the new endpoint check and extends beyond
\(41/250\). It has no local-square exception. Its checked complete
cover supplies the tiny joining interval below the rational T1
threshold; T1 supplies the remaining half-line.
The function \texttt{verify\_halfline\_gluing} checks every finite
interval join, the source cubic thresholds, and the rounding budget
using exact fractions. It explicitly asserts that the source intervals
reach the rational T1 threshold.

\paragraph{Inputs and replay.}
The analytic numerical inputs are the incoming results named
\emph{Global quadratic majorant}, \emph{Exact box bound},
\emph{Trace exclusion}, T1 (\emph{Buffered two-driver family}) and
T2 (\emph{Small buffers, one large cubic}), together with their
regenerated interval anchor data. The inherited state document's
Lemma~9 (\emph{exact convex interpolation in the buffer}) and
Theorem~18 (\emph{large buffer}) enter only through the accepted
source T1/T2 cover. The new endpoint calculation does not interpolate
in a cubic coefficient or duplicate a Gaussian buffer.

The only imported numerical helper is the incoming
\texttt{src/boxcheck.py}; no exploratory script is imported.
The new checker reads exact binary64 endpoints as rational fractions,
uses unshifted A/B elimination, and records every accepted cell.
It rejects unknown selected claim names and checks that the output
claim set equals the requested set. The outer replay regenerates the
incoming data and verifies its source claims before this extension
checker is run.

\paragraph{An immediate weak-K region using T1 alone.}
The simpler monotone transfer already gives the completely global
consequence
\[
 \boxed{\quad g\ge819/5000=0.1638
 \quad\Longrightarrow\quad
 \Lambda(G)<S+87/10^6<S+89/10^6.\quad}                         \tag{WK}
\]
This allows every \(c_1,c_2\) and every complex A/B coefficient.
For \(g\ge41/250\) use T1 directly. Otherwise apply T1 at
\(g_0=41/250\) and note that the exact rational expression
\[
 p\left((41/250)^2-(819/5000)^2\right)
 =\frac{1639}{18910000}<87/10^6
\]
is the maximum price of removing the Gaussian increment. Because T1
is a global bound, this argument has no trace-shell gap and needs no
slab-specific margin. The exact identity and all scalar inequalities
are replayed by \texttt{check\_analytic.py}.

\subsection{Coefficient disks: what complex analysis actually buys}
\label{sec:cc-poisson}
\paragraph{Provenance and status (IDE note).} From the canonical-core GPT package (20 October,
\code{conjK-canonical-core-20261020.zip}, archived in \pathref{gpt-reviews/canonical-core-20261020/}); \st{P} for the theorem, \st{C} for (CC.12).
An available interpolation tool for certified covers, \emph{not} a completed new cover. IDE checks: the Laplacian of $|w+\zeta H|$ is $|H|^2/|w+\zeta H|$,
the Cauchy--Schwarz step $\E[|H|^2/X]\ge(\E|H|)^2/\E X$, $\E|L|^2=a^2+4ab+6b^2=1181331/1562500$, $c_A=0.17932$, $c_B=0.16559$.
The Gaussian input polynomial is not holomorphic, so maximum modulus
cannot be applied to its input coordinates. However, dependence on its
\emph{coefficients} is holomorphic before taking absolute values.
This gives a legitimate Poisson comparison, provided the price is retained.

\begin{theorem}[A corrected coefficient-disk bound, \textbf{[P]}]
\label{thm:cc-poisson}
Let \(G_\zeta=G_0+\zeta H\), \(|\zeta|\le R\), with square-integrable
random variables. Suppose
\[
 \mathbb E[L\overline{G_\zeta}]=0,
 \qquad \mathbb E|G_\zeta|^2\le V_0 \quad(|\zeta|\le R).
\]
Set \(n=\mathbb E|H|^2\), \(m=\mathbb E|H|\),
\(M=\mathbb E|L|^2+V_0\), and \(c=m^2/(2\sqrt M)\).
For \(|\zeta|<R\),
\[
 \boxed{\quad
 \Lambda(G_\zeta)\le
 \int_0^{2\pi}\frac{R^2-|\zeta|^2}{|Re^{it}-\zeta|^2}
               \Lambda(G_{Re^{it}})\frac{dt}{2\pi}
 +(pn-c)(R^2-|\zeta|^2).
 \quad} \tag{CC.11}
\]
No independence between \(G_0\) and \(H\) is required.
\end{theorem}
\begin{proof}
For \(P(\zeta)=\mathbb E|L+G_\zeta|+\mathbb E|L-G_\zeta|\),
the distributional Laplacian is
\[
 \Delta P=\mathbb E\left[|H|^2
       \left(\frac1{|L+G_\zeta|}+\frac1{|L-G_\zeta|}\right)\right].
\]
The formula is understood by smoothing the norms if necessary.
Cauchy--Schwarz and the hypotheses give, for each sign,
\[
 \mathbb E\frac{|H|^2}{|L\pm G_\zeta|}
 \ge\frac{(\mathbb E|H|)^2}{\mathbb E|L\pm G_\zeta|}
 \ge\frac{m^2}{\sqrt M}.
\]
For a completely regular argument use \(\sqrt{|w|^2+\delta^2}\): its
Laplacian is at least \(|H|^2/\sqrt{|w|^2+\delta^2}\), and then let
\(\delta\downarrow0\). Thus \(P-c|\zeta|^2\) is subharmonic.
The Laplacian of \(\mathbb E|G_\zeta|^2\) is \(4n\), even when
\(G_0,H\) are correlated. Consequently
\(\Lambda(G_\zeta)+(pn-c)|\zeta|^2\) is subharmonic. Poisson's
inequality proves (CC.11).
\end{proof}

In the current remote-plus-A/B family, the orthogonality condition holds.
For a unit A direction, \(H=(|Z|^2-1)X\),
\(n=1\) and \(m=\sqrt\pi/\mathrm e\). For a unit B direction,
\(H=Z^2\bar X/\sqrt2\), \(m=\sqrt\pi/(2\sqrt2)\).
On any coefficient disk with total variance at most \(13/20\),
\(\mathbb E|L|^2=1181331/1562500\), and exact rational checks give
\[
 c_A>179/1000,\qquad c_B>165/1000. \tag{CC.12}
\]
Thus a boundary upper bound fills an A disk with error at most
\((p-179/1000)(R^2-|\zeta|^2)\), instead of the plain subharmonic
error \(p(R^2-|\zeta|^2)\). The corresponding B coefficient is
\(p-165/1000\). This reduces that interpolation allowance by at least
\(13.5\%\) and \(12.4\%\), respectively.
Applying (CC.11) successively in several coefficients gives the product
Poisson boundary average on a polydisk, with the sum of the corresponding
quadratic errors. The same variance bound must hold on the whole
polydisk, including its boundary. Thus both A directions can be treated
together without assuming independent drivers.
A sufficient disk-variance check is
\(\|G_0\|^2+2R|\mathbb E[G_0\bar H]|+R^2n\le13/20\).
The result requires a certified \emph{continuous} boundary bound;
finite phase samples alone do not suffice. It does not eliminate
phase variables or prove unrestricted A.

\paragraph{Why ordinary maximum modulus would be wrong.}
For \(p>1\), even the deterministic function
\[
 J(\zeta)=|1+\zeta|+|1-\zeta|-p|\zeta|^2
\]
has its strict global maximum at the interior point \(0\). Indeed,
\[
 |1+\zeta|+|1-\zeta|\le2\sqrt{1+|\zeta|^2}
 \le2+|\zeta|^2,
\]
so \(J(\zeta)<J(0)=2\) for \(\zeta\ne0\).
This refutes the proposed maximum-principle shortcut for penalized
norms; it is not a counterexample to K.

\subsection{Cold-start regions in the open small-buffer part}\label{sec:tdb-coldstart-regions}
\paragraph{Provenance and status (IDE note, 21 October).} A context-free AI answered the self-contained prompt \pathref{PROMPT\_GPT\_LARGE\_A\_SMALL\_BUFFER.md}
(archived in \pathref{gpt-reviews/two-driver-small-buffer-20261021/}). Its exact BigInt certificate (\code{gaussian-coupling-certificate.zip}) was only a sandbox link
and is \emph{not} available to the IDE, so the three rational certificates below are the AI's \st{C}, not replayed here. IDE checks: the conditional second moment
$\Phi$ by hand, and the global affine bounds and the three regions by an independent float quadrature of the double integral (3{,}000 random points; worst
gaps $-0.0011$, $-0.0022$, $-0.0039$ below the affine bounds; at least $0.0016$ below $2\E|L|$ on the region boundaries). Status: \st{P}+\st{C (AI)}+\st{F (IDE)}.

\paragraph{The conditioning.} For $G\in\mathcal F_2$ put $C=c_1^2+c_2^2$, $A=|\alpha|^2$, $B=|\beta|^2$, $w=g^2$, $V=C+A+B+w$, $R=|X_1|^2+|X_2|^2$ and
$\eta=\sqrt2\Re\sum_jc_j\bar\alpha_j$ (so $\eta^2\le2CA$). Given $R=r$, each $|X_j|^2$ is uniform on $[0,r]$ and the driver phases are independent, so
$\E[|H(X_j)|^2\mid r]=\frac{r^3}8-\frac{2r^2}3+r$, $\E[H(X_j)\bar X_j\mid r]=\frac{1}{\sqrt2}(\frac{r^2}3-r)$, and
\[
\Phi(t,r)=\E[|L+G|^2\mid T=t,R=r]=t(a+bt)^2+C\Bigl(\tfrac{r^3}8-\tfrac{2r^2}3+r\Bigr)+A(t-1)^2\tfrac r2+Bt^2\tfrac r4+\eta(t-1)\Bigl(\tfrac{r^2}3-r\Bigr)+w,
\]
which keeps the cubic--$A$ interference $\eta$ (the dependence lost by treating the coupling as fresh noise). Every other cross term vanishes by visible or
independent driver phase averaging. Conditional Cauchy--Schwarz gives $\Lambda(G)\le2\iint\sqrt{\Phi}\,re^{-t-r}\,dr\,dt-pV$, and
$2\sqrt\Phi\le h+\Phi/h$ with piecewise-constant rational $h$ gives affine bounds. After eliminating $\eta$ by $|\eta|\le\sqrt{2CA}$:
\[
\Lambda(G)\le\min\bigl\{1.42391-0.4495A-0.6508B+0.2216w,\ 1.46028-0.3893C-0.6676B+0.1208w,\ 1.46212-0.36(C+A)-0.67328B+0.08554w\bigr\}
\]
for every $g\ge0$.

\begin{theorem}[\st{P}+\st{C (AI)}]\label{tdb-thm:coldstart}
For $G\in\mathcal F_2$ with $0\le g\le0.164$ and all complex phases, each of $A+B\ge0.116+0.494g^2$, or $C\ge0.236$, or $V\ge0.285$ implies
$\Lambda(G)<2\E|L|=\frac{387}{500}\sqrt\pi\le S$. These hold without the $\hat c$ caps or the per-driver inequalities.
\end{theorem}
This improves the energy exclusion $|\alpha|^2+|\beta|^2\ge4/25$ to about $0.13$, and the variance exclusion $V\ge13/20$ to $V\ge0.285$, on this family. It also
covers $g<0.1$, where Theorem~\ref{tdb-thm:B} gave no cubic threshold, whenever $c_1^2+c_2^2\ge0.236$.

\paragraph{Spherical decomposition (same AI, \st{P}).} Writing $(X_1,X_2)=\sqrt R\,\omega$ with $\omega$ uniform on $S^3$,
$H(\sqrt r\omega_j)=\frac{r^{3/2}}{\sqrt2}h_j(\omega)+\sqrt{2r}(\frac r3-1)\omega_j$ with harmonic $h_j=\omega_j(|\omega_j|^2-\frac23)$. The cubic's linear spherical
component merges with the $A$-coupling before the absolute value is taken. The zonal-kernel $L^1$ norm is \emph{not} a valid bound even for $Y_1+Y_2$:
exactly $\|Y_1+Y_2\|_1/\|Y_1+Y_2\|_2>\|Y_1\|_1/\|Y_1\|_2$. The projection constants of Defant et al.\ are $C(\mathbb S)$ operator norms, not $L^2\to L^1$ bounds.
The rest of the open region is unresolved.

\paragraph{Two further cold-start answers (IDE-replayed, 21 October).} Two more context-free AIs answered the same prompt; their packages
\code{two-driver-partial-proof.zip} and \code{two-driver-certificate.zip} are archived in \pathref{gpt-reviews/two-driver-small-buffer-20261021/}. Both Node certificates
(\code{node:assert} only, BigInt outward intervals, exact tail moments) PASS in the IDE: 3.4\,s and 18\,s, the latter over about $3.1\cdot10^5$ cells. Both condition on
$(T,|X_1|^2,|X_2|^2)$ and retain the cubic--$A$ cross term $K=\sqrt2\sum_jc_j\Re\alpha_j$. They prove the global affine bounds
\[
\begin{aligned}
\Lambda&\le1.4215-0.4399A-0.6430B+0.2466w, &\Lambda&\le1.4430-0.3650C-0.0336A-0.6456B+0.2219w,\\
\Lambda&\le1.4476-.3917C-.6518B+.1894w, &\Lambda&\le1.4204-.4382A-.6442B+.2505w,\\
\Lambda&\le1.3981-.6374B+.4800w, &\Lambda&\le1.4540-.3714(C+A)-.6617B+.1310w,
\end{aligned}
\]
with benchmarks $S\ge1.3739269$ (Jensen cells) and $S\ge1.373145$. IDE float check: all six bounds lie at least $0.0035$ above
$2\E\sqrt{\E[|Y|^2\mid T,S_1,S_2]}-pV$ at 1{,}500 random points of $\mathcal F_2$. \textbf{Combined regions} (with $g\le0.164$, all phases): $\Lambda<S$ whenever
$C\ge0.2$, or $A+B\ge0.124$, or $A+1.46B\ge0.1233$, or $B\ge0.06$, or $C+A+B\ge0.228$, or $V\ge0.285$, or $\min$ of the affine bounds $\le1.3739$. The remaining open
part of the small-buffer two-driver region lies inside $C<0.2$, $A+B<0.116$, $B<0.06$, $C+A+B<0.228$. The second AI also reports a normalisation erratum in
Defant et al., Thms~5.2--5.3 (denominator $2^{n-1+|p-q|/2}$, not $2^{n+|p-q|/2}$), and the harmonic split $Q=\mathcal H+(R^2/4-1)\kappa$,
$\kappa=\sum_j(\sqrt2c_j+\alpha_j)X_j$, which shows the $A$-coupling entering the radial component.


\paragraph{The coupling-free slice, ideal target (fourth cold-start AI, IDE-replayed, 21 October).} A fourth context-free AI answered the same prompt
(\code{two\_driver\_slice\_proof.zip}, archived in \pathref{gpt-reviews/two-driver-small-buffer-20261021/slice-proof/}). For $\alpha=\beta=0$, $G$ is
independent of $Z$ and rotation invariant. From $\E|W|=\int_0^\infty(1-\E J_0(k|W|))k^{-2}\,dk$ in $\mathbb R^2$ it gets
\[
\Lambda=2\int_0^\infty\bigl(1-\Xi(k)\Psi(c_1k)\Psi(c_2k)e^{-g^2k^2/4}\bigr)\frac{dk}{k^2}-p(c_1^2+c_2^2+g^2),\qquad
\Xi(k)=\E J_0(k|L|),\ \Psi(t)=\E J_0(t|H|).
\]
\textbf{Theorem \st{P}+\st{C (AI)}.} If $\alpha=\beta=0$, $g\le0.164$, $\max(c_1,c_2)\le\hat c(g)$ for $g\ge0.1$ and $c_1^2+c_2^2+g^2\le0.65$, then
$\Lambda\le S-10^{-4}$. This sharpens fact~6 of the prompt ($S+10^{-5}$) to the ideal target on this slice. \emph{Ingredients} (arb balls, python-flint):
$S\ge s(0.040345)\ge1.3739535256614$; $\Xi$ at the 4\,200 nodes of composite 3-point Gauss--Legendre on $[0,70]$ ($h=1/20$); $\Psi,\Psi'$ at $t=j/100$,
$j\le5651$, with cubic Hermite interpolation ($|\Psi''''|\le\E|H|^4=44$) and a degree-7 Taylor series for $t\le0.01$; the GL3 error through
$|((1-\Phi)/k^2)^{(6)}|\le\E|L+G|^8/56$ (Minkowski); the tail $|\Xi(k)|\le\|\Delta^3f_L\|_1k^{-6}\le29860.39\,k^{-6}$. Convexity of $2\E|L+G|$ in $(c_1,c_2,g)$
bounds each box by the trilinear interpolant of its vertices: 323 boxes, minimum gap $1.02\cdot10^{-4}$, at $g\uparrow0.164$, $c_1=c_2\approx0.026$.
\emph{IDE replay:} both covering parts PASS (20\,s and 5\,s); \code{s\_lower.py} and \code{lap\_rig.py} reproduce $S_{\rm low}$ and $29860.3807$; 50 random entries of the $\Xi$, $\Psi$ and $\Psi'$ tables
agree with independent mpmath quadrature to $5\cdot10^{-15}$. The IDE checked by hand the radial formula, the Hermite, Taylor and GL3 remainders,
$\E|H|^4=44$, and $\Delta=4\partial_u(u\,\partial_u)$ for the density $f_L=e^{-T}/(\pi u'(T))$.
\emph{Not proved:} shared couplings break the factorisation. There each driver has a three-variable characteristic function over 11 parameters. The AI's float
evidence: near the slice maximum the second variation in $(\alpha,\beta)$ is negative definite (about $-0.43|\alpha|^2-0.94|\beta|^2+0.29|\alpha||\beta|$ at $g=0.164$). Shared $A$-couplings gain
only where $c\ge0.2$, which is already $\ge8\cdot10^{-3}$ below $S$. Sixty random coupled points give $\Lambda-S\le-7.2\cdot10^{-3}$.
